\documentclass[12pt]{amsart}
\usepackage{ amsmath, amsthm, amsfonts, amssymb, color}
 \usepackage{mathrsfs}
\usepackage{amsfonts, amsmath}
 \usepackage{amsmath,amstext,amsthm,amssymb,amsxtra}
 \usepackage{txfonts} %also pxfonts
 \usepackage[colorlinks, citecolor=blue,pagebackref,hypertexnames=false]{hyperref}
 \allowdisplaybreaks
 \usepackage{pgf,tikz}
 \usepackage{multirow}%newtablepackage1
 \usepackage{diagbox} %newtablepackage2
\usepackage{cite}
 \usepackage{tikz}

\usetikzlibrary{decorations.pathreplacing}

 %\usepackage{showkeys}
 \textwidth =166mm
 \textheight =232mm
\marginparsep=0cm
\oddsidemargin=2mm
\evensidemargin=2mm
\headheight=13pt
\headsep=0.8cm
\parskip=0pt
\hfuzz=6pt
\widowpenalty=10000
 \setlength{\topmargin}{-0.2cm}

\DeclareMathOperator*{\essinf}{ess\ inf}
\DeclareMathOperator*{\esssup}{ess\ sup}
\setcounter {tocdepth}{1}
\begin{document}

 \baselineskip 16.6pt
\hfuzz=6pt

\widowpenalty=10000

\newtheorem{cl}{Claim}
\newtheorem{theorem}{Theorem}[section]
\newtheorem{proposition}[theorem]{Proposition}
\newtheorem{coro}[theorem]{Corollary}
\newtheorem{lemma}[theorem]{Lemma}
\newtheorem{definition}[theorem]{Definition}
\newtheorem{nota}[theorem]{Notation}
%\theoremstyle{definition}
\newtheorem{assum}{Assumption}[section]
\newtheorem{example}[theorem]{Example}
\newtheorem{remark}[theorem]{Remark}
\renewcommand{\theequation}
{\thesection.\arabic{equation}}

\def\SL{\sqrt H}

\newcommand{\mar}[1]{{\marginpar{\sffamily{\scriptsize
        #1}}}}

\newcommand{\as}[1]{{\mar{AS:#1}}}
\newcommand{\cE}{\mathcal{E}}
\newcommand\R{\mathbb{R}}
\newcommand\RR{\mathbb{R}}
\newcommand\CC{\mathbb{C}}
\newcommand\D {{\rm Dunkl}}
\newcommand\NN{\mathbb{N}}
\newcommand\N {{\bf N}}
\newcommand\ZZ{\mathbb{Z}}
\newcommand\HH{\mathbb{H}}
\newcommand\Z{\mathbb{Z}}
\def\RN {\mathbb{R}^n}
\renewcommand\Re{\operatorname{Re}}
\renewcommand\Im{\operatorname{Im}}

\newcommand{\mc}{\mathcal}

\def\hs{\hspace{0.33cm}}
\newcommand{\la}{\alpha}
\def \l {\alpha}
\newcommand{\eps}{\tau}
\newcommand{\pl}{\partial}
\newcommand{\supp}{{\rm supp}{\hspace{.05cm}}}
\newcommand{\x}{\times}
\newcommand{\lag}{\langle}
\newcommand{\rag}{\rangle}

\newcommand\wrt{\,{\rm d}}

\newcommand{\norm}[2]{|#1|_{#2}}
\newcommand{\Norm}[2]{\|#1\|_{#2}}

\title[Schatten class properties of Dunkl--Riesz transform commutators]{Weyl--Chamber geometry, Poincar\'e inequality meets Schatten class properties of Dunkl--Riesz transform commutators}



\author{Zhijie Fan}
\address{Zhijie Fan, School of Mathematical Sciences, South China Normal University, Guangzhou 510631, China}
\email{fanzhj3@mail2.sysu.edu.cn}







\author{Ji Li}
\address{Ji Li, Department of Mathematics, Macquarie University,NSW2109, Australia}
\email{ji.li@mq.edu.au}



\author{Dmitriy Zanin}
\address{Dmitriy Zanin, University of New South Wales, Kensington, NSW 2052, Australia}
\email{d.zanin@unsw.edu.au}






  \date{\today}

 \subjclass[2010]{47B10, 42B20, 43A85}
\keywords{Schatten class, Riesz transform commutators, Dunkl operator, Besov space, Nearly weakly orthogonal sequence}




\begin{abstract}

\end{abstract}

\maketitle


\tableofcontents

%\newpage
\section{Introduction}
\begin{definition}{\rm\label{beso}
Suppose $1< p< \infty$. We say that a  function $f\in L^p_{{\rm loc}}(\mathbb{R}^{N},dw)$ belongs to {Dunkl} Besov space $B^{p}_{\D}(\mathbb{R}^N)$ if
\begin{align*}
\|f\|_{B^{p}_{\D}(\mathbb{R}^N)}:=\left(\int_{\mathbb{R}^{N}}\int_{\mathbb{R}^{N}}\left(\frac{d(x,y)}{\|x-y\|}\right)^{p}\frac{|f(x)-f(y)|^p}{w(B(x,d(x,y)))^2}dw(x)dw(y)\right)^{1/p}<+\infty.
\end{align*}}
\end{definition}

\subsection{Notation}
Conventionally, we let $\mathbb{N}$ denote the set of positive integers and set $\mathbb{Z}_+ := \{0\} \cup \mathbb{N}$. Throughout the paper, the symbols $C$ and $c$ denote positive constants (possibly varying from line to line) that may depend only on $N$, $p$, and the structural data $(R,k)$ of the Dunkl setting, but are independent of the main parameters (e.g., the symbol $b$ and auxiliary decay factors).
 We write $A \lesssim B$ to mean that $A \leq C B$ for some constant $C > 0$, and $A \sim B$ to mean that both $A \lesssim B$ and $B \lesssim A$ hold.
For \(\lambda>0\) and a cube \(Q\) with sidelength \(\ell(Q)\), write \(\lambda Q\) for the cube concentric with \(Q\) and sidelength \(\lambda\,\ell(Q)\). Similarly, for every ball \(B=B(x,r)\), set \(\lambda B:=B(x,\lambda r)\). For any $1\leq p\leq\infty$, we denote by $p'$ the conjugate of $p$, which satisfies $\frac{1}{p}+\frac{1}{p'}=1$. We denote the average of a function $f$ over a $\mu$-measurable set $E$ by
\begin{align}\label{defaverage}
(f)_E:=\fint_{E}f(x)d\mu(x):=\frac{1}{\mu(E)}\int_{E}f(x)d\mu(x).
\end{align}
For a subset $E \subset \mathbb{R}^{N}$, we denote by $\chi_E$ its characteristic function.

Denote by $\{e_j\}_{j=1}^N$ the standard orthonormal basis on $\mathbb{R}^N$.
%For two measurable Borel sets $E$ and $F$, we use the notation ${\rm dist}(E,\,F)$ to denote the distance between them.


\section{Preliminaries}
\setcounter{equation}{0}
\subsection{Harmonic analysis in the Dunkl setting}
Let $\mathbb{R}^N$ be the Euclidean space equipped with the usual scalar product $\langle \cdot,\,\cdot\rangle $ and the induced norm $\|\cdot\|$. For a nonzero vector $v\in \mathbb{R}^N$, the reflection $\sigma_v$ with respect to the hyperplane $v^{\perp}$ orthogonal to $v$ is defined by
\begin{align*}
	\sigma_v x:= x-2\frac{\langle x,v\rangle}{\|v\|^2}v.
\end{align*}
A finite set $R\subset \mathbb{R}^N\backslash \{0\}$ is called a root system if $\sigma_v(R)=R$ for every $v\in R$. A root system $R$ is called a normalized reduced root system if  $\|v\|^2=2$ for every $v\in R$.
The finite group $G$ generated by the reflections $\{\sigma_v\}_{v\in R}$ is called the Weyl group (reflection group) of the root system $R$.
Let $k: R \to [0,\infty)$ be a multiplicity function, which is a $G$-invariant function  fixed throughout the paper. We will denote by $\mathcal O(x)=\{\sigma(x): \sigma \in G\}$ the $G$-orbit of a point $x\in\mathbb{R}^N$, and $\mathcal O(E)=\bigcup_{x \in E} \mathcal O(x)$ the $G$-orbit of $E\subset \mathbb{R}^N$.

Define the Dunkl measure $dw$ on $\mathbb{R}^N$ by
$$
dw(x) :=\prod_{v \in R} |\langle x,v\rangle|^{k(v)} dx,
$$
where $dx$ stands for the Lebesgue measure on $\mathbb{R}^N$. Let $B(x,r):=\{y\in \mathbb{R}^N:\|x-y\|< r\}$ be the Euclidean ball centred at $x$ with radius $r>0$. Let $\N:=N+\sum_{v\in R}k(v)$. Note that
$$w(B(tx,tr))=t^{\N}w(B(x,r)),\ {\rm for}\ x\in\mathbb{R}^N,\ t,r>0.$$
Furthermore,
\begin{align}\label{twosideinequ}
  w(B(x,r))\sim r^N \prod_{v \in R}(|\langle x, v \rangle |+r)^{k(v)},
\end{align}
where the implicit constant only depends on $N,R$ and $k$.  It follows from \eqref{twosideinequ} that there exists a constant $C\geq 1$ such that for every $x\in \mathbb{R}^N$,
\begin{align}\label{doublinginequality}
C^{-1}\left(\frac{r_2}{r_1}\right)^N\leq \frac{w(B(x,r_2))}{w(B(x,r_1))}\leq C\left(\frac{r_2}{r_1}\right)^\N,\ {\rm for}\ 0<r_1<r_2,
\end{align}
so the numbers $\N$ and $N$ are called the upper and lower homogeneous dimensions, respectively.

The Dunkl operators $T_\xi$ are defined by
$$T_\xi f(x):=\partial_\xi f(x)+\sum_{v\in R}\frac{k(v)}{2}\langle v,\xi\rangle\frac{f(x)-f(\sigma_v(x))}{\langle v,x\rangle},$$
where $\partial_\xi f$ denotes the directional derivative of $f$ in the direction $\xi$. Furthermore, the Dunkl Laplacian is defined by $$\Delta_{\D}f(x):=\sum_{j=1}^N T_{e_j}^2f(x)=\Delta_{\mathbb{R}^N}f(x)+\sum_{v\in R}k(v)\delta_v f(x),$$
where $\Delta_{\mathbb{R}^N}$ is the Laplacian on $\mathbb{R}^N$ and $\delta_v$ is defined by
$$\delta_v f(x):=\frac{\partial_v f(x)}{\langle v,x\rangle}-\frac{f(x)-f(\sigma_v(x))}{\langle v,x\rangle^2}.$$
In this article, we consider the Dunkl-Riesz transforms defined by $R_{\D,j}:=T_{e_j}(-\Delta_{\D})^{-1/2}$, $j=1,2,\ldots,N$.


\begin{definition}{\rm
The Weyl chambers are defined as the closures of connected components of the set
  \begin{align}
    \left\{x\in\mathbb{R}^N:\langle x,v\rangle \neq 0\ \ { for} \ { all} \ v\in R\right\}.\end{align}}
\end{definition}
Define $
	d(x,y):=\min_{\sigma\in G}\|x-\sigma(y)\|
$
as the distance between two $G$-orbits $\mathcal O(x)$ and $\mathcal O(y)$. Recall from \cite[Chapter VII, proof of Theorem 2.12]{MR514561} that
\begin{align}\label{distanceequa}
d(x,y)=\|x-y\|\ {\rm if}\ {\rm and}\ {\rm only}\ {\rm if}\ x,y\in\Gamma
\end{align}
for some closed Weyl chamber $\Gamma$.

It is clear that $\mathcal O(B(x,r)) = \{y\in \mathbb{R}^N: d(x,y)<r\}$. Moreover,
\[
w(B(x,r)) \leq w(\mathcal O(B(x,r))) \leq |G|w(B(x,r))
\]
for all $x \in \mathbb{R}^N$ and $r>0$. We also denote $B^d(x,r):= \{y\in \mathbb{R}^N: d(x,y)<r\}$, i.e., $B^d(x,r)=\mathcal O(B(x,r))$.

Denote by $K_{\D,j}(x,y)$ the integral kernel of $R_{\D,j}$. We recall the kernel estimates for the Riesz transforms in the Dunkl setting proved in \cite{MR4733961}. Although these estimates may not be sharp (see the comment following Theorem~2.1 in \cite[p.~10]{MR4733961}), they suffice for our purposes.
\begin{lemma}\label{kernelestimate}
There exists a constant $C$ such that for $j\in\{1,2,\ldots,N\}$ and every $x,y$ with $d(x,y)\neq 0$,
\begin{align*}
|K_{\D,j}(x,y)|\leq C\frac{d(x,y)}{\|x-y\|}\frac{1}{w(B(x,d(x,y)))}.
\end{align*}
\end{lemma}
\begin{proof}
The proof is given in \cite[Theorem 1.1]{MR4733961}.
\end{proof}

\begin{lemma}\label{nondegenerate}
For $j\in\{1,2,\ldots,N\}$, and for every ball $B=B(x_0,r)\subset \mathbb{R}^N$, there is another ball $\widehat{B}=B(y_0,r)$ with $y_0=x_0+5re_j$ such that for every $(x,y)\in B\times \widehat{B}$, $K_{\D,j}(x,y)$ does not change sign and satisfies
\begin{align*}
|K_{\D,j}(x,y)|\geq\frac{C}{w(B(x_0,r))}.
\end{align*}
\end{lemma}
\begin{proof}
The lower bound is given in \cite[Theorem 1.2]{MR4733961}, while the fact that $K_{\D,j}(x,y)$ does not change sign can be seen in the proof of \cite[Theorem 1.2]{MR4733961}. Indeed, with the choice $y_0:=x_0+5re_j$, we have $y_j-x_j\geq 3r$ for all $(x,y)\in B\times \widehat{B}$. Together with the fact that the heat kernel $h_t(x,y)$ associated with the Dunkl Laplacian is  positive, the representation
$$K_{\D,j}(x,y)=-\frac{C}{2}\int_0^\infty \frac{y_j-x_j}{t}h_t(x,y)\frac{dt}{\sqrt{t}},$$
implies that $K_{\D,j}(x,y)<0$ (and thus, does not change sign) on $ B\times \widehat{B}$.
\end{proof}



\subsection{Systems of dyadic cubes}


\begin{definition}\label{Defsystem}{\rm
Let $X\subseteq \mathbb{R}^N$ be a nonempty open set. A countable family
$
    \mathscr{D}
    := \bigcup_{k\in\mathbb{Z}}\mathscr{D}_{k}, \
    \mathscr{D}_{k}
    :=\{Q^k_{\alpha}\colon \alpha\in \mathscr{A}_k\},
$
of Borel sets $Q^k_{\alpha}\subseteq X$ is called \textit{a
system of dyadic cubes over  $X$}  if it has the following properties:

\begin{enumerate}
  \item[(i)] For each $k\in\mathbb{Z}$, we have $$    X
    = \bigcup_{\alpha\in \mathscr{A}_k} Q^k_{\alpha};$$\smallskip

  \item[(ii)] $\text{If }\ell\geq k\text{, then either }
        Q^{\ell}_{\beta}\subseteq Q^k_{\alpha}\text{ or }
        Q^k_{\alpha}\cap Q^{\ell}_{\beta}=\varnothing$;\smallskip

  \item[(iii)] $
    \text{For each }(k,\alpha)\text{ and each } \ell\leq k,
    \text{ there exists a unique } \beta
    \text{ such that }Q^{k}_{\alpha}\subseteq Q^\ell_{\beta};
$

\smallskip

\item[(iv)] For each $(k,\alpha)$ there exist $2^N$ indices $\beta$ such that
    $$ Q^{k+1}_{\beta}\subseteq Q^k_{\alpha},\ {\rm and}\
        Q^k_{\alpha} =\bigcup_{{\substack{Q\in\mathscr{D}_{k+1}\\ Q\subseteq Q^k_{\alpha}}}}Q;$$

\smallskip

\item[(v)] For each $(k,\alpha)$, one has $$B_X(x^k_{\alpha},A_1 2^{-k})
    \subseteq Q^k_{\alpha}\subseteq B_X(x^k_{\alpha},A_2 2^{-k})
    =: B(Q^k_{\alpha})$$
    for some positive constants $A_1,A_2$ depending only on the geometry of $X$, {where we set $B_X(x,r):=B(x,r)\cap X$ for all $x\in X$ and $r>0$.}

\smallskip

\item [(vi)]  If $\ell\geq k$ and
   $Q^{\ell}_{\beta}\subseteq Q^k_{\alpha}$, then
   $$B(Q^{\ell}_{\beta})\subseteq B(Q^k_{\alpha}).$$
\end{enumerate}



   }
\end{definition}

Each $Q^k_{\alpha}$ is called a \textit{dyadic cube of generation} $k$, with center $x^k_{\alpha}\in Q^k_{\alpha}$ and side-length $\ell(Q^k_{\alpha})=2^{-k}$.
\begin{nota}
For every dyadic cube $Q\in\mathscr{D}$, we denote by $c_Q$ and $\ell(Q)$ its center and side-length, respectively.
\end{nota}

\begin{nota}\label{standardsystem}
Let $$\mathbb{D}^0:=\bigcup_{k\in\mathbb{Z}}\mathbb{D}_k^0$$ be the standard system of dyadic cubes over $\mathbb{R}^N$.
Here each dyadic cube is a product of half-open intervals
\[
[m_1 2^{-k}, (m_1+1)2^{-k}) \times \cdots \times [m_N 2^{-k}, (m_N+1)2^{-k}),
\]
with $(m_1,\ldots,m_N)\in \mathbb{Z}^N$. Moreover, we let $$\mathbb{D}_{\Gamma}:=\bigcup_{k\in\mathbb{Z}}\mathbb{D}_{\Gamma,k}$$ be the system of dyadic cubes over a fixed chamber $\Gamma$, whose  construction will be given in Corollary \ref{Construction over gamma}.
\end{nota}


\begin{lemma}\label{Nsystems}
There exist $2^N$ systems of dyadic cubes $\mathbb{D}^1,\mathbb{D}^2,\ldots,\mathbb{D}^{2^N}$ such that for any ball $B=B(x_B,r_B)\subset \mathbb{R}^N$ with $\frac{2^{-k-2}}{3}\leq r_B< \frac{2^{-k-1}}{3}$ for some $k\in\mathbb{Z}$,  there exists a cube ${\displaystyle Q\in\bigcup_{\nu=1}^{2^N}\mathbb{D}^\nu_k}$ such that
$$B\subset Q\subset 12\sqrt{N} B.$$
\end{lemma}

\begin{proof}
{\bf $\bullet$ Step 1.}  {\it Recall the one-dimensional setting given in \cite[Section 2]{MR3420475}.}

Let $\mathbb{D}^{(0)}(\mathbb{R})$ be the standard system of dyadic intervals over $\mathbb{R}$ (see Notation \ref{standardsystem}).
Let $\delta=1/3$ and define
$$
   \mathbb D^{(\delta)}(\mathbb R)=\bigcup_{k\in\mathbb Z} \mathbb D^{(\delta)}_k(\mathbb R),
$$
where for $k\geq0$,
$$
    \mathbb D^{(\delta)}_k(\mathbb R)=\{I+\delta \mid I\in \mathbb D_k^{(0)}(\mathbb R)\},
$$
while for $k<0$ and $k$ even,
$$
    \mathbb D^{(\delta)}_k(\mathbb R)
    =\Big\{\Big[{m\over 2^k}+\delta+\sum_{j=(k/2)+1}^0 2^{-2j},\ {m+1\over 2^k}
    +\delta+\sum_{j=(k/2)+1}^0 2^{-2j}\Big)\ \Big|\ m\in
    \mathbb Z\Big\}.
$$
These choices together with the nestedness property property completely
determine the collections $\mathbb D^{(\delta)}_k(\mathbb R)$ for $k<0$, $k$ odd.

Fix an interval $J\subset\mathbb R$ with $2^{-k-1}/3\le |J|<2^{-k}/3$ for some $k\in\mathbb Z$. Now we set
$A_k=\{ m\cdot 2^{-k} \mid m\in\mathbb Z \}$ for every fixed $k$ and
\begin{enumerate}
    \item[(1)]$A_k^{(\delta)}=\{ \delta+ m\cdot 2^{-k} \mid m\in\mathbb Z
        \}$ for $k\geq 0$,

    \item[(2)]$A_k^{(\delta)}=\{ \delta+
        \sum_{j=(k/2)+1}^{0}2^{-2j}+ m\cdot 2^{-k} \mid m\in\mathbb Z
        \}$ for $k<0$, $k$ even, and

    \item[(3)]$A_k^{(\delta)}=\{ \delta+
        \sum_{j=(k-1)/2+1}^{0}2^{-2j}+ m\cdot 2^{-k} \mid
        m\in\mathbb Z \}$ for $k<0$, $k$ odd.
\end{enumerate}
Note that for any two different points $a,b\in A_k\cup
A_k^{(\delta)}$, we have $|a-b|\geq 2^{-k}/3>|J|$. Thus,
there is at most one element of $A_k\cup A_k^{(\delta)}$ contained in
$J$. Hence, $J\cap A_k=\emptyset $ or $J\cap
A_k^{(\delta)}=\emptyset $. Therefore, $J$ must be contained in
some dyadic interval $I\in \mathbb D_k^{(0)}$ or $I\in \mathbb D^{(\delta)}_k$ and
$|I|=2^{-k}\leq 6 |J|$.

{\bf $\bullet$ Step 2.} {\it We extend the construction to $\mathbb{R}^N$ ($N\geq 1$).}

For $\vartheta=(\vartheta_1,\vartheta_2,\ldots,\vartheta_N)\in\{0,\delta\}^N$, define
$$\mathbb{D}^{(\vartheta)}:=\bigcup_{k\in\mathbb{Z}}\mathbb{D}_k^{(\vartheta)},\quad \mathbb{D}_k^{(\vartheta)}:=\left\{\prod_{j=1}^N I_j: I_j\in \mathbb{D}_k^{(\vartheta_j)}\right\}.$$
Applying Step 1 to $\pi_i(B)$, where $\pi_i(B)$ denotes the $i$-th coordinate projection, we obtain $I_i\in\mathbb{D}_k^{(0)}$ or $I_i\in\mathbb{D}_k^{(\delta)}$ with $\pi_i(B)\subset I_i$ and $|I_i|=2^{-k}\leq 12 r_B$. Let $\vartheta_i\in\{0,\delta\}$
be the corresponding shift so that $I_i\in\mathbb D_k^{(\vartheta_i)}$, set $\vartheta:=(\vartheta_1,\ldots,\vartheta_N)$, and let $Q=\prod_{j=1}^NI_j\in\mathbb{D}_k^{(\vartheta)}$. It follows that $B\subset Q$ and $\ell(Q)=2^{-k}\leq 12r_B$. Hence, for every $y\in Q$,
$$\|y-x_B\|\leq \sqrt{N}\ell(Q)\leq 12\sqrt{N}r_B.$$
This implies that $Q\subset B(x_B,12\sqrt{N}r_B)$. Enumerating $\big\{\mathbb D^{(\vartheta)}:\vartheta\in\{0,\delta\}^N\big\}$ as
$\mathbb D^1,\ldots,\mathbb D^{2^N}$, we complete the proof of Lemma \ref{Nsystems}.
\end{proof}



%Hence, for any ball $B\subset \mathbb{R}^N$, via reducing to each axis in $\mathbb{R}^N$
%and repeating the argument in $\mathbb{R}$, we see that there exist $k\in\mathbb Z$ and a cube ${\displaystyle Q\in\bigcup_{\nu=1}^{2^N}\mathbb{D}^\nu_k}$ such that
%$$B\subset Q\subset \kappa_N B,$$
%where $\kappa_N$ is a positive constant depending only on $N$.


\subsection{Haar Basis}\label{HaarSection}



Given a system of dyadic cubes $\mathscr{D}$ over $\mathbb{R}^N$, we now recall the standard way to construct an orthonormal Haar system corresponding to $\mathscr{D}$. For each $k\in\mathbb{Z}$ and $Q\in\mathscr{D}_k$, denote by $Ch(Q):=\{R\in\mathscr{D}_{k+1}:R\subseteq Q\}$. Moreover,  we let
$$h_Q^0:=w(Q)^{-1/2}\chi_Q.$$
Fix an enumeration $\{Q_i\}_{i=1}^{2^N}$ of $Ch(Q)$. Let
$$\hat{Q}_i:=\bigcup_{j=i}^{2^N}Q_j$$
and
$$h_Q^i:=\sqrt{\frac{w(Q_i)w(\hat{Q}_{i+1})}{w(\hat{Q}_i)}}\Bigg(\frac{\chi_{Q_i}}{w(Q_i)}-\frac{\chi_{\hat{Q}_{i+1}}}{w(\hat{Q}_{i+1})}\Bigg),\quad i=1,2,\ldots,2^N-1.$$


\begin{lemma}\label{Haarexpansion}
For any $1<p<\infty$ and $f\in
    L^p(\mathbb{R}^N,dw)$, we have
    \[
        f(x)
        =  \sum_{Q\in\mathscr{D}}\sum_{i=1}^{2^N-1}
            \langle f,h^{i}_{Q}\rangle_{L^2(\mathbb{R}^N,dw)} h^{i}_{Q}(x), %\quad \text{if } \mu(X)=\infty,
    \]
    where the sum converges both in the
    $L^p(\mathbb{R}^N,dw)$ norm and pointwise almost everywhere.
 \end{lemma}


\begin{lemma}
The family of Haar functions $\{h_{Q}^{i}\}_{Q\in\mathscr{D},0\leq i\leq 2^N-1}$ has the following properties:
    \begin{itemize}
\item[(i)] $h_{Q}^{i}$ is supported on $Q$ for $i\in\{0,1,\ldots,2^N-1\}$;
\item[(ii)] $h_{Q}^{i}$ is constant on each $R\in Ch(Q)$ for $i\in\{0,1,\ldots,2^N-1\}$;
\item[(iii)] $\int_{\mathbb{R}^N} h_{Q}^{i}(x)dw(x)= 0$, for $i\in\{1,2,\ldots,2^N-1\}$;
\item[(iv)] $\langle h_{Q}^{i},h_{Q}^{j}\rangle_{L^2(\mathbb{R}^N,dw)} = 0$ for
$i \neq j$, $i$, $j\in\{1, 2,\ldots, 2^N- 1\}$;
\item[(v)] The collection
$\{h_{Q}^{i}\}_{i = 0}^{2^N-1}$
is an orthogonal basis for the vector space $V(Q)$ of all functions on $Q$ that are constant on each sub-cube $R\in Ch(Q)$.
\item[(vi)]
            $
 \Norm{h_{Q}^{i}}{L^p(\mathbb{R}^N,dw)}
    \approx w(Q)^{\frac{1}{p} - \frac{1}{2}}$
   \quad for~1 $\leq p \leq \infty$ and $i\in\{1,2,\ldots,2^N-1\}$;

    \end{itemize}
\end{lemma}


\subsection{Nearly weakly  orthonormal sequences}
Rochberg and Semmes \cite{RS} introduced the notion of \emph{nearly weakly  orthonormal (NWO)} sequences of functions. This notion has since proved highly effective in characterizing Schatten class membership of  singular integral commutators in a variety of contexts (see \cite{FLLXarxiv,HKarxiv,MR4756021,MR4552581,MR4873796,LLW,RS,Zhangwei1,Zhangwei2}). The properties of NWO sequences were originally investigated in the Euclidean setting (equipped with Lebesgue measure), but the extension to spaces of homogeneous type is straightforward once a system of dyadic cubes is available in this general framework (see \cite[Section 13]{HKarxiv} for a complete proof).  For our purposes, we adapt  the extension to the Dunkl setting $(\mathbb{R}^N,\|\cdot\|,w)$ rather than recalling the general  definition
of NWO sequences and the general extension to spaces of homogeneous type.

{\it Let $\mathscr{D}$ be a system of dyadic cubes. A sequence of functions $\{e_Q\}_{Q\in\mathscr{D}}$ is an NWO sequence if
\begin{enumerate}
  \item $\supp e_Q\subseteq cQ$ for some positive constant $c$ independent of $Q$
  \item $\|e_Q\|_{L^s(\mathbb{R}^{N},dw)}\leq w(Q)^{1/s-1/2}$ for some $s>2$.
\end{enumerate}
}
\begin{lemma}\label{NWOLEMMAA}
Let $1<p<\infty$, and let $\mathscr{D}$ be a
system of dyadic cubes over $\mathbb{R}^N$. For any compact operator $ T $ on $ L ^2 (\mathbb R ^{N},dw)$, we have
\begin{equation*}
\Bigl(
\sum_{Q\in \mathscr{D}} |\langle T e_Q, f_Q \rangle_{L^2(\mathbb{R}^N,dw)} | ^{p}
\Bigr) ^{1/p} \leq C_{N,p,w} \| T \| _{S ^{p}},
\end{equation*}
where $\{e_Q\}_{Q\in\mathscr{D}}$ and $\{f_Q\}_{Q\in\mathscr{D}}$ are NWO sequences.
\end{lemma}
\begin{proof}
This property can be found in  \cite[(1.10)]{RS} (see also \cite[Section 7]{HKarxiv}).
\end{proof}

\begin{lemma}\label{NWOLEMMAA2}
Let $0<p<\infty$, and let $\mathscr{D}$ be a
system of dyadic cubes over $\mathbb{R}^N$. Assume that
$$A=\sum_{Q\in\mathscr{D}}\lambda_Q\langle \cdot,\,e_Q\rangle_{L^2(\mathbb{R}^N,dw)} f_Q,$$
where $\{e_Q\}_{Q\in\mathscr{D}}$ and $\{f_Q\}_{Q\in\mathscr{D}}$ are NWO sequences and $\{\lambda_Q\}_{Q\in\mathscr{D}}$ is a sequence of scalars. We have
\begin{align}\label{nwose}
\|A\|_{S^{p}}\leq C_{N,p,w} \|\{\lambda_Q\}_Q\|_{\ell^{p}}.
\end{align}
\end{lemma}
\begin{proof}
This property can be found in  \cite[(1.9)]{RS} (see also \cite[Section 7]{HKarxiv}).
\end{proof}

\section{Analytic tools for Weyl chamber geometry}
\setcounter{equation}{0}

{\bf Convention. }
Throughout this paper, we will use the following notation:
\begin{itemize}
  \item Let $R
  \subset \mathbb{R}^{N}$ be a finite normalized reduced root system;
  \item Let
$W:=\operatorname{span}(R)$ with $N_{0}:=\dim W$;
\item Let
$
  \Omega:=\mathbb{R}_+^{N_{0}}\times \mathbb{R}^{N-N_{0}};
$
\item Let $\Gamma_+$ be a fixed open Weyl chamber (see Definition \ref{def:weyl_chamber_positive_roots}) and $\Gamma$ be its closure.
\end{itemize}


\begin{definition}{\rm \cite[Definition ~2.7]{HeckmanCoxeterGroups}}\label{def:weyl_chamber_positive_roots} A connected component of \begin{align}
    \left\{x\in\mathbb{R}^N:\langle x,v\rangle \neq 0\ \ { for} \ { all} \ v\in R\right\}
  \end{align} is called an open
Weyl chamber. Its closure is called a Weyl chamber.  Define the corresponding sets of \emph{positive} and \emph{negative} roots by
\[
R_{+}:=\{\alpha\in R:\ \langle x,\alpha\rangle>0\ \text{for all }x\in \Gamma_{+}\},
\qquad
R_{-}:=-R_{+}.
\]
\end{definition}

\begin{definition}{\rm \cite[Definition~2.9]{HeckmanCoxeterGroups}}\label{def:simple_roots}
A root $\alpha\in R_{+}$
is called simple in $R_{+}$ if it cannot be written as
\[
\alpha=\lambda_{1}\alpha_{1}+\lambda_{2}\alpha_{2}
\quad\text{with }\lambda_{1},\lambda_{2}\ge 1\ \text{and}\ \alpha_{1},\alpha_{2}\in R_{+}.
\]
The set $\Delta\subset R_{+}$ of all simple roots is called the
simple system.
\end{definition}

\begin{lemma}\label{lem:simple_system_associated_to_chamber}
With the notation given in Definition \ref{def:simple_roots},  $\Delta$ is a basis of $W$. If we denote by
$
\Delta=\{\alpha_{1},\alpha_{2},\ldots,\alpha_{N_{0}}\},
$
then the open and closed chambers can be written as
\begin{align}
\Gamma_{+}
&=\Bigl\{x\in\mathbb{R}^{N}:\ \langle x,\alpha_{i}\rangle>0,\ i=1,\ldots,N_{0}\Bigr\},
\label{eq:open_chamber_by_simple_roots}\\
\Gamma
&=\Bigl\{x\in\mathbb{R}^{N}:\ \langle x,\alpha_{i}\rangle\ge 0,\ i=1,\ldots,N_{0}\Bigr\},
\label{eq:closed_chamber_by_simple_roots}
\end{align}
and every root $\alpha\in R$ admits a unique expansion
\[
\alpha=\sum_{i=1}^{N_{0}}c_{i}\alpha_{i}
\]
such that either $c_{i}\ge 0$ for all $i$ (when $\alpha\in R_{+}$), or $c_{i}\le 0$ for all $i$
(when $\alpha\in R_{-}$). In particular, the simple system $\Delta$ is uniquely determined by $\Gamma_{+}$
(equivalently, by $\Gamma$).
\end{lemma}

\begin{proof}
%We regard $R$ as a normalized root system in the Euclidean space $W$ endowed with the
%restricted inner product. Since $\sigma_{\alpha}(\alpha)=-\alpha$ and $\sigma_{\alpha}(R)=R$
%for every $\alpha\in R$, we have $R=-R$.

Let $\Delta=\{\alpha_{1},\ldots,\alpha_{n}\}\subset R_{+}$ be the set of simple roots in $R_{+}$. By \cite[Proposition~2.10]{HeckmanCoxeterGroups}, every $\alpha\in R_{+}$
is a nonnegative linear combination of $\alpha_{1},\ldots,\alpha_{n}$. Since $R=R_{+}\cup R_{-}$ (this follows from the fact that for each $\alpha\in R$, the continuous function
$x\mapsto\langle x,\alpha\rangle$ does not vanish on the connected set $\Gamma_{+}$) and $R_{-}=-R_{+}$, it follows that
$R\subset\operatorname{span}(\Delta)$ and hence,
\[
W=\operatorname{span}(R)\subset\operatorname{span}(\Delta).
\]
The reverse inclusion $\operatorname{span}(\Delta)\subset W$ is immediate from $\Delta\subset R\subset W$,
so $\operatorname{span}(\Delta)=W$. On the other hand, it follows from \cite[Theorem~2.16]{HeckmanCoxeterGroups} that the simple roots are linearly
independent. Consequently, $\Delta$ is a basis
of $W$, and therefore $n=\dim W=N_{0}$. Moreover, every $\alpha\in R$ has a unique expansion
$\alpha=\sum_{i=1}^{N_{0}}c_{i}\alpha_{i}$. If $\alpha\in R_{+}$, the coefficients can be chosen
nonnegative by \cite[Proposition~2.10]{HeckmanCoxeterGroups}. If $\alpha\in R_{-}=-R_{+}$, then
$-\alpha\in R_{+}$ has an expansion with all coefficients nonnegative, and multiplying by $-1$ gives
an expansion of $\alpha$ with all coefficients nonpositive.
The descriptions \eqref{eq:open_chamber_by_simple_roots} and \eqref{eq:closed_chamber_by_simple_roots} follow from \cite[Corollary~2.11]{HeckmanCoxeterGroups}.


To see that $\Delta$ is uniquely determined by $\Gamma_+$, note that $\Gamma_+$ determines $R_{+}$
by Definition~\ref{def:weyl_chamber_positive_roots}, and the set of simple roots in $R_{+}$ is
characterized intrinsically by Definition~\ref{def:simple_roots}. Finally, since $\Gamma_+$ is the (unique) interior component of $\Gamma$, specifying $\Gamma$ determines $\Gamma_+$ and vice versa. This completes the proof of Lemma \ref{lem:simple_system_associated_to_chamber}.
\end{proof}
\begin{nota}
With the notation given in Lemma \ref{lem:simple_system_associated_to_chamber}, denote by
 $$\Delta=\{\alpha_{1},\ldots,\alpha_{N_{0}}\}$$ the simple system associated with $\Gamma$.
\end{nota}

\begin{lemma}\label{chambergeometry}
Each Weyl chamber $\Gamma$ admits the decomposition
\[
\Gamma = K(\Gamma)\times W^\perp,
\qquad
K(\Gamma):=\{w\in W:\ \langle w,\alpha\rangle\ge 0\ \text{for all }\alpha\in\Delta\}.
\]
\end{lemma}

\begin{proof}
Write $x=w+z$ with $w\in W$ and $z\in W^{\perp}$. Since $\alpha_{i}\in W$, we have $\langle x,\alpha_{i}\rangle=\langle w,\alpha_{i}\rangle$ for all
$i=1,\ldots,N_0$. This, together with \eqref{eq:closed_chamber_by_simple_roots}, implies that $x\in \Gamma$ if and only if
$
w\in K(\Gamma).
$
This implies that
$
\Gamma
=
K(\Gamma)\times W^{\perp},
$
and thus completes the proof.
\end{proof}

%
%\begin{coro}
%Let $R\subset\mathbb{R}^N$ be a finite normalized reduced root system, let $W:=\operatorname{span}(R)$ and $N_{0}:=\dim W$, and let $D$ be any Weyl chamber in $\mathbb{R}^N$. Then there exists a linear isomorphism
%\[
%\Phi:\mathbb{R}^N\longrightarrow \mathbb{R}^N
%\]
%such that
%\[
%\Phi(D)=\mathbb{R}_+^{\,N_{0}}\times \mathbb{R}^{\,N-N_{0}}.
%\]
%\end{coro}
%
%\begin{proof}
%By the theorem \ref{chambergeometry}, we have the product decomposition
%\[
%D=K(D)\times W^\perp,
%\qquad
%K(D)=\{w\in W:\ \langle w,\alpha\rangle\ge0\ \ \forall\,\alpha\in\Delta\},
%\]
%where $\Delta=\Delta(D)=\{\alpha_1,\dots,\alpha_{N_{0}}\}\subset R$ consists of $N_{0}$ linearly independent facet normals; in particular, $K(D)$ is a simplicial cone in $W$ given by the intersection of exactly $N_{0}:=\dim W$ closed half-spaces with linearly independent outward normals $\Delta$.
%
%Define $L:W\to\mathbb{R}^{N_{0}}$ by
%\[
%L(w):=\big(\langle w,\alpha_1\rangle,\dots,\langle w,\alpha_{N_{0}}\rangle\big).
%\]
%Relative to the basis $\Delta$ of $W$, the matrix of $L$ is the Gram matrix
%$G=(\langle \alpha_j,\alpha_i\rangle)_{i,j}$, which is nonsingular because
%$\Delta$ is linearly independent. Hence $L$ is a linear isomorphism $W\!\xrightarrow{\;\sim\;}\mathbb{R}^{N_{0}}$. From the description of $K(D)$ it follows that
%\[
%L\big(K(D)\big)=\{t\in\mathbb{R}^{N_{0}}:\ t_i\ge0\ \text{for all }i\}=\mathbb{R}_+^{\,N_{0}}.
%\]
%Choose any linear isomorphism $U:W^\perp\!\xrightarrow{\;\sim\;}\mathbb{R}^{N-N_{0}}$, and define
%\[
%\Phi:\mathbb{R}^N=W\oplus W^\perp\longrightarrow \mathbb{R}^{N_{0}}\oplus \mathbb{R}^{N-N_{0}}\cong\mathbb{R}^N,
%\qquad
%\Phi(w+z):=\big(L(w),\,U(z)\big).
%\]
%Then $\Phi$ is a linear isomorphism and
%\[
%\Phi(D)=\Phi\big(K(D)\times W^\perp\big)=L\big(K(D)\big)\times U\big(W^\perp\big)
%=\mathbb{R}_+^{\,N_{0}}\times \mathbb{R}^{\,N-N_{0}},
%\]
%as claimed.
%\end{proof}



\begin{coro}\label{coro:chamber-orthant}
Let $\Phi_{W^{\perp}}$ be any linear isomorphism from $W^\perp$ to $\mathbb{R}^{N-N_{0}}$, and let
\[
\Phi_W:W\longrightarrow\mathbb{R}^{N_{0}},
\qquad
\Phi_W(w):=\big(\langle w,\alpha_1\rangle,\dots,\langle w,\alpha_{N_{0}}\rangle\big),
\]
\[
\Phi:\mathbb{R}^N=W\oplus W^\perp\longrightarrow
\mathbb{R}^{N_{0}}\oplus\mathbb{R}^{N-N_{0}}\cong\mathbb{R}^N,
\qquad
\Phi(w+z):=\big(\Phi_W(w),\,\Phi_{W^\perp}(z)\big).
\]
Then $\Phi$ is a linear isomorphism and
\[
\Phi(\Gamma)=\mathbb{R}_+^{N_{0}}\times \mathbb{R}^{N-N_{0}}.
\]
\end{coro}

\begin{proof}
By Lemma~\ref{chambergeometry} we have $\Gamma=K(\Gamma)\times W^\perp$, where
\[
K(\Gamma)=\{w\in W:\ \langle w,\alpha_i\rangle\ge0,\ i=1,\dots,N_{0}\}.
\]
Since $\Delta$ is a basis of $W$, the Gram matrix
$G=(\langle \alpha_i,\alpha_j\rangle)_{i,j=1}^{N_{0}}$ is nonsingular. Hence,
$\Phi_W:W\longrightarrow\mathbb{R}^{N_{0}}$ is a linear isomorphism. Moreover,
for $w\in W$ we have $w\in K(\Gamma)$ if and only if $\Phi_{W}(w)\in\mathbb{R}_+^{N_{0}}$,
so
\[
\Phi_{W}\big(K(\Gamma)\big)=\mathbb{R}_+^{N_{0}}.
\]
By construction, $\Phi$ is a linear isomorphism (as a direct sum of isomorphisms),
and therefore
\[
\Phi(\Gamma)
=\Phi\big(K(\Gamma)\times W^\perp\big)
=\Phi_{W}\big(K(\Gamma)\big)\times \Phi_{W^{\perp}}\big(W^\perp\big)
=\mathbb{R}_+^{N_{0}}\times \mathbb{R}^{N-N_{0}},
\]
which proves the desired conclusion.
\end{proof}


\begin{coro}\label{Construction over gamma}
 Let $\Gamma$ be an arbitrary Weyl chamber. There exists a system of dyadic cubes over $\Gamma$.
\end{coro}\begin{proof}
Let $\Phi:\mathbb{R}^N\to\mathbb{R}^N$ be the linear isomorphism from Corollary~\ref{coro:chamber-orthant} so that
\[
\Phi(\Gamma)=\Omega:=\mathbb{R}_+^{N_0}\times \mathbb{R}^{N-N_0}.
\]

For each $k\in\mathbb{Z}$, define $\mathbb{D}^0_{\Omega,k}$ to be the family of half-open cubes
\[
Q=\prod_{i=1}^{N} I_i,
\qquad
I_i=\big[m_i2^{-k},(m_i+1)2^{-k}\big),
\]
where $m_i\in\mathbb{Z}_+$ for $1\le i\le N_0$ and $m_i\in\mathbb{Z}$ for $N_0<i\le N$. Set
\[
\mathbb{D}^0_\Omega:=\bigcup_{k\in\mathbb{Z}}\mathbb{D}^0_{\Omega,k}.
\]
It is standard that $\mathbb{D}^0_\Omega$ is a system of dyadic cubes over $\Omega$.

Next, for each $k\in\mathbb{Z}$, define
\[
\mathbb{D}_{\Gamma,k}:=\big\{\Phi^{-1}(Q): Q\in \mathbb{D}^0_{\Omega,k}\big\},
\qquad
\mathbb{D}_{\Gamma}:=\bigcup_{k\in\mathbb{Z}}\mathbb{D}_{\Gamma,k}.
\]
Since $\Phi:\mathbb{R}^N\to\mathbb{R}^N$ is a linear isomorphism, we see that $\mathbb{D}_{\Gamma}$ is a system of dyadic cubes over $\Gamma$.
\end{proof}




\begin{nota}
 Let $\Phi$ be the linear isomorphism given in Corollary~\ref{coro:chamber-orthant}. For every Borel set $E\subset\Omega$, define
\begin{align}\label{defmu}
  \mu(E):=w\big(\Phi^{-1}(E)\big).
\end{align}
\end{nota}
\begin{lemma}\label{mulemma}
We have
  \begin{equation}\label{pushforward-density}
d\mu(u,z)=\phi(u)\,du\,dz,\quad \phi(u)=|{\rm det}(\Phi)|^{-1}\prod_{v\in R_+}\Big(\sum_{i=1}^{N_{0}} a_{v,i}\,u_i\Big)^{2k(v)}.
\end{equation}
where for each $v\in R_+$, the coefficients $a_{v,i}\ge0$ are uniquely determined by the simple-root expansion
\(v=\sum_{i=1}^{N_0} a_{v,i}\alpha_i\) and satisfy $(a_{v,1},\dots,a_{v,N_0})\neq\mathbf 0$.
\end{lemma}
\begin{proof}
Since $\Phi(\Gamma)=\Omega$, we have $\Phi^{-1}(\Omega)=\Gamma$. Hence, for every nonnegative Borel function
$f$ on $\Omega$,
\[
\int_{\Omega} f(u,z)\,d\mu(u,z)
=\int_{\Gamma} f\big(\Phi(x)\big)\,dw(x).
\]
We make the linear change of variables $(u,z)=\Phi(x)$, so that $x=\Phi^{-1}(u,z)$. Since $\Phi$ is linear and invertible, the Jacobian is constant and
\[
dx = |{\rm det}(\Phi)|^{-1}\,du\,dz.
\]
It follows from the definition of $\Phi$ that
$
u_i=\langle x,\alpha_i\rangle$ for $i=1,\dots,N_0.
$
This implies that
\begin{align}\label{changevariable}
\langle x,v\rangle
=\Big\langle x,\sum_{i=1}^{N_0} a_{v,i}\alpha_i\Big\rangle
=\sum_{i=1}^{N_0} a_{v,i}\langle x,\alpha_i\rangle
=\sum_{i=1}^{N_0} a_{v,i}u_i.
\end{align}
Substituting these identities into the previous integral, we obtain
\begin{align*}
\int_{\Omega} f(u,z)\,d\mu(u,z)
&=\int_{\Gamma} f\big(\Phi(x)\big)\Big(\prod_{v\in R_+}|\langle x,v\rangle|^{2k(v)}\Big)\,dx \\
&=\int_{\Omega} f(u,z)\,
|{\rm det}(\Phi)|^{-1}\prod_{v\in R_+}\Big(\sum_{i=1}^{N_0} a_{v,i}u_i\Big)^{2k(v)}\,du\,dz.
\end{align*}
Therefore, \eqref{pushforward-density} holds. This completes the proof of Lemma~\ref{mulemma}.
\end{proof}


\begin{lemma}\label{lem:gallery-simple}
Let $\Gamma_1$ and $\Gamma_2$ be two Weyl chambers. Then there exist an integer $M\ge 0$ and
$\alpha_1,\dots,\alpha_M\in R$ such that, if we set $\Gamma^{(0)}:=\Gamma_2$ and
\[
\Gamma^{(k)}:=\sigma_{\alpha_k}\big(\Gamma^{(k-1)}\big),\qquad k=1,\dots,M,
\]
then the following hold:
\begin{enumerate}
  \item\label{geo1} $\alpha_k\in \Delta(\Gamma^{(k-1)})$ for each $k=1,\dots,M$;
  \item\label{geo2} $\Gamma^{(M)}=\Gamma_1$.
\end{enumerate}
\end{lemma}
\begin{proof}
If $\Gamma_1=\Gamma_2$, we take $M=0$ and there is nothing to prove. Assume $\Gamma_1\neq \Gamma_2$.
Since the Weyl group $G$ acts transitively on the set of Weyl chambers (see \cite[Theorem 2.12 (2)]{HeckmanCoxeterGroups}), there exists
$g\in G$ such that $\Gamma_1=g(\Gamma_2)$.

Let $\Delta(\Gamma_2)=\{\beta_1,\beta_2,\ldots,\beta_{N_0}\}$ be the simple system associated with $\Gamma_2$.
By \cite[Theorem 2.12(4)]{HeckmanCoxeterGroups}, the Weyl group $G$ is generated by the simple reflections $\sigma_{\beta_1},\sigma_{\beta_2},\ldots,\sigma_{\beta_{N_0}}$. Hence, there exist indices
$i_1,\ldots,i_M\in\{1,\ldots,N_0\}$ such that
\[
g=\sigma_{\beta_{i_1}}\sigma_{\beta_{i_2}}\cdots \sigma_{\beta_{i_M}}.
\]
Define
\[
w_0:=\mathrm{Id},\qquad w_k:=\sigma_{\beta_{i_1}}\cdots \sigma_{\beta_{i_k}}\quad (k=1,\ldots,M),
\]
and set
\[
\Gamma^{(k)}:=w_k(\Gamma_2)\qquad (k=0,1,\ldots,M).
\]
It follows that $\Gamma^{(0)}=\Gamma_2$ and $\Gamma^{(M)}=w_M(\Gamma_2)=g(\Gamma_2)=\Gamma_1$, which implies the statement \eqref{geo2}.

For each $k=1,\ldots,M$, define
\[
\alpha_k:=w_{k-1}(\beta_{i_k})\in R.
\]
The rest of the proof is divided into two steps.

{\bf $\bullet$ Step 1.} {\it Verify that $\Gamma^{(k)}=\sigma_{\alpha_k}(\Gamma^{(k-1)})$.}

Since each $w\in G$ is orthogonal, for any root $\beta\in R$ one has the conjugation identity
\[
w\,\sigma_{\beta}\,w^{-1}=\sigma_{w(\beta)},\qquad w\in G.
\]
Applying this with $w=w_{k-1}$ and $\beta=\beta_{i_k}$ yields
\[
\sigma_{\alpha_k}=\sigma_{w_{k-1}(\beta_{i_k})}=w_{k-1}\sigma_{\beta_{i_k}}w_{k-1}^{-1}.
\]
Therefore,
\[
\sigma_{\alpha_k}(\Gamma^{(k-1)})
= w_{k-1}\sigma_{\beta_{i_k}}w_{k-1}^{-1}\big(w_{k-1}(\Gamma_2)\big)
= w_{k-1}\sigma_{\beta_{i_k}}(\Gamma_2)
= w_k(\Gamma_2)
= \Gamma^{(k)}.
\]

{\bf $\bullet$ Step 2.} {\it Verify that $\alpha_k\in\Delta(\Gamma^{(k-1)})$.}

Recall that if $\Gamma$ is a Weyl chamber associated with simple system $\Delta(\Gamma)$, then
\[
\Gamma=\{x\in\mathbb R^N:\ \langle x,\gamma\rangle\ge 0\ \text{for all }\gamma\in\Delta(\Gamma)\}.
\]
Applying $w_{k-1}\in G$ to the above description for $\Gamma_2$ and using
$\langle w_{k-1}x,\,\beta\rangle=\langle x,\,w_{k-1}^{-1}\beta\rangle$, we obtain
\[
\Gamma^{(k-1)}=w_{k-1}(\Gamma_2)
=\Bigl\{x\in\mathbb R^N:\ \langle x,\,w_{k-1}(\beta_j)\rangle\ge 0\ \text{for }j=1,\ldots,N_0\Bigr\}.
\]
By the uniqueness of the simple system associated with a Weyl chamber
(see Lemma~\ref{lem:simple_system_associated_to_chamber}), we conclude that
\[
\Delta(\Gamma^{(k-1)})=\Delta\big(w_{k-1}(\Gamma_2)\big)=w_{k-1}\big(\Delta(\Gamma_2)\big)=\{w_{k-1}(\beta_1),\ldots,w_{k-1}(\beta_{N_0})\}.
\]
In particular, $\alpha_k=w_{k-1}(\beta_{i_k})\in\Delta(\Gamma^{(k-1)})$, proving the first statement \eqref{geo1}.

This completes the proof of Lemma \ref{lem:gallery-simple}.
\end{proof}


\begin{definition}
We say that $R'\subset\Omega$ is an axis--parallel rectangle if
\[
R' := I_1\times\cdots\times I_N,
\]
with $I_i=[I_i^-,I_i^+]\subset[0,\infty)$ for $1\le i\le N_{0}$ and $I_i=[I_i^-,I_i^+]\subset\mathbb{R}$ for $i>N_{0}$.
\end{definition}

Denote by
$$B_\Omega(x_0,r):=\{x\in\Omega:\ \|x-x_0\|<r\}.$$
The following lemma demonstrates that $(\Omega,\|\cdot\|,d\mu)$ is a doubling metric measure space.
\begin{lemma}\label{lem:ball-2sided-N}
There exists a constant $C\ge 1$ such that
for every $x_0=(u_0,z_0)\in\Omega$ and every $r>0$,
\begin{equation}\label{eq:ball-2sided-Omega}
C^{-1}\,r^N\prod_{v\in R_+}\Big(L_{v,0}+r\Big)^{2k(v)}
\ \le\
\mu\big(B_\Omega(x_0,r)\big)
\ \le\
C\,r^N\prod_{v\in R_+}\Big(L_{v,0}+r\Big)^{2k(v)},
\end{equation}
where
\[
L_v(u):=\sum_{i=1}^{N_{0}} a_{v,i}u_i,
\qquad
L_{v,0}:=L_v(u_0)=\sum_{i=1}^{N_{0}} a_{v,i}(u_0)_i.
\]
In particular, there exists a constant $C_{\mathrm{dbl},\Omega}\geq 1$ such that for every $x\in\Omega$ and every $r>0$,
\[
\mu\big(B_\Omega(x,2r)\big)\ \le\ C_{\mathrm{dbl},\Omega}\;\mu\big(B_\Omega(x,r)\big).
\]
\end{lemma}

\begin{proof}
Fix $x_0=(u_0,z_0)\in\Omega$ and $r>0$. For each $v\in R_+$ define the linear form
\[
L_v(u):=a_v\cdot u,\qquad a_v:=(a_{v,1},\dots,a_{v,N_{0}})\in \mathbb{R}_+^{N_{0}}\setminus\{0\}.
\]
Since $R_+$ is finite and each $a_v\neq 0$, we have
\[
L_*:=\max_{v\in R_+}\|a_v\|<\infty,
\qquad
m_*:=\min_{v\in R_+}\sum_{i=1}^{N_{0}} a_{v,i}>0.
\]

\noindent
\textbf{Upper bound.}
Let $x=(u,z)\in B_\Omega(x_0,r)$. Then
\[
\|u-u_0\|\ \le\ \|x-x_0\|\ \le\ r.
\]
Hence, for every $v\in R_+$,
\[
|L_v(u)-L_{v,0}|
=\big|a_v\cdot (u-u_0)\big|
\le \|a_v\|\cdot\|u-u_0\|
\le L_*\,r.
\]
Since $L_v(u)\ge 0$ on $\Omega$ (all coefficients and coordinates are nonnegative), we obtain
\[
L_v(u)\le L_{v,0}+L_*r\le C_1\,(L_{v,0}+r),
\qquad
C_1:=\max\{1,L_*\}.
\]
Therefore,
\[
\phi(u)
=\prod_{v\in R_+} L_v(u)^{2k(v)}
\le C_1^{2\sum_{v\in R_+}k(v)}\prod_{v\in R_+}(L_{v,0}+r)^{2k(v)}
=:C_2\prod_{v\in R_+}(L_{v,0}+r)^{2k(v)}.
\]
Integrating over $B_\Omega(x_0,r)$ and using
$|B_\Omega(x_0,r)|\le |B(x_0,r)|=\omega_N r^N$, where $\omega_N$ is the volume of the unit ball in $\mathbb{R}^N$, gives
\[
\mu\big(B_\Omega(x_0,r)\big)
=\int_{B_\Omega(x_0,r)}\phi(u)\,dx
\le C_2\,|B(x_0,r)|\,\prod_{v\in R_+}(L_{v,0}+r)^{2k(v)}
\le C_3\,r^N\prod_{v\in R_+}(L_{v,0}+r)^{2k(v)},
\]
with $C_3:=C_2\,\omega_N$.

\medskip\noindent
\textbf{Lower bound.}
We construct a smaller Euclidean ball $B_c$ contained in $B_\Omega(x_0,r)$ on which $\phi$ admits a uniform
positive lower bound in terms of $\prod_{v\in R_+}(L_{v,0}+r)^{2k(v)}$. To this end, we first choose
\[
\lambda:=1+\frac{L_*}{m_*}>1,
\qquad
\eta:=\frac{1}{4\big(\sqrt{N_{0}}\,\lambda+1\big)},
\qquad
\theta:=\lambda\,\eta.
\]
These parameters depend only on $L_*$ and $m_*$, and satisfy
\[
\theta> \eta>0,
\qquad
\sqrt{N_{0}}\,\theta+\eta=\eta(\sqrt{N_{0}}\lambda+1)=\frac14<1.
\]
Let $\mathbf{1}_{N_{0}}:=(1,\dots,1)\in\mathbb{R}^{N_{0}}$ and define the center
\[
x_c:=(u_c,z_c):=\big(u_0+\theta r\,\mathbf{1}_{N_{0}},\ z_0\big),
\qquad
B_c:=B(x_c,\eta r).
\]
We {\it claim} that $B_c\subset B_\Omega(x_0,r)$. To see this inclusion, we first apply the triangle inequality to see that for any $x\in B_c$,
\[
\|x-x_0\|\le \|x-x_c\|+\|x_c-x_0\|
\le \eta r+\theta r\sqrt{N_{0}}
\le r,
\]
since $\sqrt{N_{0}}\theta+\eta<1$. Hence $B_c\subset B(x_0,r)$. Moreover, for any $x=(u,z)\in B_c$, we have $\|u-u_c\|\le \|x-x_c\|\le \eta r$. This implies that for each $1\le i\le N_{0}$,
\[
u_i\ge (u_c)_i-\eta r=(u_0)_i+(\theta-\eta)r\ge 0,
\]
since $u_0\in\mathbb{R}_+^{N_{0}}$ and $\theta\ge\eta$. Thus $u\in\mathbb{R}_+^{N_{0}}$ and $x\in\Omega$. Therefore,  $B_c\subset B_\Omega(x_0,r)$. This finishes  the proof of the  claim.

\smallskip
Next we obtain a uniform lower bound for $L_v(u)$ on $B_c$.
Write any $x=(u,z)\in B_c$ as
\[
u=u_0+\theta r\,\mathbf{1}_{N_{0}}+\Delta_u,
\qquad
z=z_0+\Delta_z,
\qquad
\|\Delta_u\|^2+\|\Delta_z\|^2\le (\eta r)^2,
\]
so in particular $\|\Delta_u\|\le \eta r$.
Then for all $x=(u,z)\in B_c$ and all $v\in R_+$,
\begin{align*}
L_v(u)
&=a_v\cdot u
=a_v\cdot u_0+\theta r\,a_v\cdot\mathbf{1}_{N_{0}}+a_v\cdot\Delta_u \\
&\ge L_{v,0}+\theta r\sum_{i=1}^{N_{0}} a_{v,i}-\|a_v\|\,\|\Delta_u\|
\ge L_{v,0}+rK_0\geq c_0(L_{v,0}+r).
\end{align*}
with
\[
K_0:=\theta m_*-\eta L_*=\eta(\lambda m_*-L_*)=\eta m_*>0,\qquad c_0:=\min\{1,K_0\}\in(0,1].
\]
Consequently, for all $u\in B_c$,
\[
\phi(u)
=\prod_{v\in R_+}L_v(u)^{2k(v)}
\ge c_0^{2\sum_{v\in R_+}k(v)}\prod_{v\in R_+}(L_{v,0}+r)^{2k(v)}
=:C_4\prod_{v\in R_+}(L_{v,0}+r)^{2k(v)}.
\]
This, in combination with the fact that $B_c\subset B_\Omega(x_0,r)$, yields
\[
\mu\big(B_\Omega(x_0,r)\big)
\ge \mu(B_c)
=\int_{B_c}\phi(u)\,dx
\ge C_4\,|B_c|\prod_{v\in R_+}(L_{v,0}+r)^{2k(v)}
= C_5\,r^N\prod_{v\in R_+}(L_{v,0}+r)^{2k(v)},
\]
where $C_5:=C_4\,\omega_N\eta^N>0$.

The second statement follows immediately from the first one. Indeed, by the first statement,
\[
\mu\big(B_\Omega(x_0,2r)\big)
\le C_3\,(2r)^N\prod_{v\in R_+}(L_{v,0}+2r)^{2k(v)}
\le C_3\,2^{N+\mathcal{K}}\,r^N\prod_{v\in R_+}(L_{v,0}+r)^{2k(v)}
\le C_3C_5^{-1}\,2^{N+\mathcal{K}}\,\mu\big(B_\Omega(x_0,r)\big),
\]
where we set $\mathcal{K}:=2\sum_{v\in R_+}k(v)$. Therefore, $(\Omega,\|\cdot\|,\mu)$ is doubling with the doubling constant $C_{\mathrm{dbl},\Omega}:=C_3C_5^{-1}\,2^{N+\mathcal{K}}$.
 This finishes the proof of Lemma \ref{lem:ball-2sided-N}.
\end{proof}
Let $\Phi:\mathbb{R}^N\longrightarrow\mathbb{R}^N$ be the linear isomorphism given in  Corollary~\ref{coro:chamber-orthant} such that
$
\Phi(\Gamma)=\Omega.
$
Since $\Phi$ is linear and invertible, there exist positive constants $c_\Phi$ and $C_\Phi$ such that for any $x,y\in\mathbb{R}^N$,
\begin{equation}\label{eq:biLip-N}
c_\Phi\,\|x-y\|\ \le\ \|\Phi(x)-\Phi(y)\|\ \le\ C_\Phi\,\|x-y\|.
\end{equation}
Denote by $B_\Gamma(x,r):=\{y\in\Gamma:\ \|y-x\|<r\}$. It follows from \eqref{eq:biLip-N} that for any $x_0\in\Gamma$ and $r>0$,
\begin{align}\label{beforedoublingball}
B_{\Omega}(\Phi(x_0),c_\Phi r)\subset \Phi(B_{\Gamma}(x_0,r))\subset B_{\Omega}(\Phi(x_0),C_\Phi r).
\end{align}
This, together with Lemma  \ref{lem:ball-2sided-N}, yields
\begin{align}\label{doublingball}
  \mu(\Phi(B_{\Gamma}(x_0,r))) \simeq \mu(B_{\Omega}(\Phi(x_0),C_\Phi r)).
\end{align}
\begin{lemma}\label{construction-rectangle}
  For all $x_0\in\Gamma$ and $r>0$, one can find an axis--parallel rectangle $R'$ on $\Omega$ such that
\begin{align}\label{step3goal}
B_{\Omega}(\Phi(x_0),C_\Phi r)\subset R'\subset B_{\Omega}(\Phi(x_0),\sqrt{N}C_\Phi r).
\end{align}
Furthermore, we have
\begin{align}\label{Rdoubling}
  \mu(R')\simeq \mu(B_{\Omega}(\Phi(x_0),C_\Phi r)).
\end{align}
%\begin{align}\label{step3goal}
%\Phi(B_\Gamma(x_0,r))\subset R'\quad \text{and}\quad
%\mu(\Phi(B_\Gamma(x_0,r)))\simeq \mu(R').
%\end{align}
\end{lemma}
\begin{proof}
We write
$
(u_0,z_0):=\Phi(x_0)\in\Omega.
$
and define the axis--parallel rectangle
\begin{equation}\label{eq:rectangle-choice}
R':=\Big(\prod_{i=1}^{N_{0}} \big[\max\{0,(u_0)_i-C_\Phi r\},\ (u_0)_i+C_\Phi r\big]\Big)
\times
\Big(\prod_{j=1}^{N-N_{0}}\big[(z_0)_j-C_\Phi r,\ (z_0)_j+C_\Phi r\big]\Big).
\end{equation}

On the one hand, it follows immediately that $B_\Omega(\Phi(x_0),C_\Phi r)\subset R'$. On the other hand, since $R'$ is a rectangle with half-side lengths bounded by $C_\Phi r$ in each coordinate, we have
\begin{align*}
R'\subset B_\Omega(\Phi(x_0),\sqrt{N}\,C_\Phi r).
\end{align*}
The second statement follows from the first one together with Lemma \ref{lem:ball-2sided-N}. The proof of Lemma \ref{construction-rectangle} is finished.
\end{proof}
The following theorem demonstrates that $(\Gamma,\|\cdot\|,dw)$ is a doubling metric measure space.


\begin{theorem}\label{main:dunkldoubling}
There exists a constant $C\ge 1$ such that
for every $x\in\Gamma$ and every $r>0$,
\begin{equation}\label{eq:ball-2sided-Gamma}
C^{-1}\,r^N \prod_{v \in R_+}(|\langle x, v \rangle |+r)^{2k(v)}
\ \le\
w\big(B_\Gamma(x,r)\big)
\ \le\
C\,r^N \prod_{v \in R_+}(|\langle x, v \rangle |+r)^{2k(v)}.
\end{equation}
In particular, there exists a constant $C_{\mathrm{dbl},\Gamma}\geq 1$ such that for every $x\in\Gamma$ and every $r>0$,
\[
w\big(B_\Gamma(x,2r)\big)\ \le\ C_{\mathrm{dbl},\Gamma}\;w\big(B_\Gamma(x,r)\big).
\]
Moreover, for all $x\in\Gamma$ and $r>0$, we have
\begin{align}\label{globallocalvol}
w(B_\Gamma(x,r))\sim w(B(x,r)).
\end{align}
\end{theorem}
\begin{proof}
The main strategy of the proof is to transfer the ball estimates from $(\Omega,\|\cdot\|,\mu)$ back to $(\Gamma,\|\cdot\|,dw)$. Fix $x\in\Gamma$ and $r>0$. It follows from the equality \eqref{defmu} that
\begin{equation}\label{eq:pushforward-ball}
w\big(B_\Gamma(x,r)\big)=\mu\big(\Phi(B_\Gamma(x,r))\big).
\end{equation}
Applying $\mu$ to the inclusions \eqref{beforedoublingball} and using \eqref{eq:pushforward-ball} yields
\begin{equation}\label{eq:wball-between-muballs}
\mu\big(B_{\Omega}(\Phi(x),c_\Phi r)\big)
\ \le\
w\big(B_\Gamma(x,r)\big)
\ \le\
\mu\big(B_{\Omega}(\Phi(x),C_\Phi r)\big).
\end{equation}
Write $\Phi(x)=(u,z)\in\Omega$. For each $v\in R_+$ recall the linear form
$L_{v,0}:=\sum_{i=1}^{N_0}a_{v,i}u_i$.
Combining the inequality \eqref{eq:wball-between-muballs} with Lemma~\ref{lem:ball-2sided-N} and the equality \eqref{changevariable}, we obtain that
\begin{equation}\label{eq:wball-in-terms-of-L}
w\big(B_\Gamma(x,r)\big)\ \simeq\ r^N\prod_{v\in R_+}\big(L_{v,0}+r\big)^{2k(v)}=r^N\prod_{v\in R_+}\big(|\langle x,v\rangle|+r\big)^{2k(v)}.
\end{equation}
This proves \eqref{eq:ball-2sided-Gamma}. The second statement follows immediately from \eqref{eq:ball-2sided-Gamma}. The third statement follows from the inequalities \eqref{eq:ball-2sided-Gamma} and \eqref{twosideinequ}. This completes the proof of Theorem~\ref{main:dunkldoubling}.
\end{proof}

\section{Test Section}
\begin{lemma}\label{lem:Dunkl-fractional-kernel-derivative}
Let $K_{\D}(x,y)$ denote the integral kernel of
$(-\Delta_{\D})^{-1/2}$. Then, for every pair of multi-indices
$\alpha,\beta$, there exists a constant $C_{\alpha,\beta}>0$ such that
\begin{equation}\label{eq:Dunkl-fractional-kernel-derivative-lemma}
\left|
\partial_x^\alpha\partial_y^\beta K_{\D}(x,y)
\right|
\leq
C_{\alpha,\beta}
d(x,y)^{-(\N-1+|\alpha|+|\beta|)}
\end{equation}
for all $x,y\in\mathbb R^N$ satisfying $d(x,y)>0$.
\end{lemma}

\begin{proof}
Let $h_t(x,y)$ be the heat kernel of $e^{t\Delta_{\D}}$. By
\cite[Theorem~4.1(d)]{MR4014802}, for every pair of multi-indices
$\alpha,\beta$, there exist constants
$C_{\alpha,\beta},c_{\alpha,\beta}>0$ such that
\begin{equation}\label{eq:Dunkl-heat-mixed-lemma}
\left|
\partial_x^\alpha\partial_y^\beta h_t(x,y)
\right|
\leq
C_{\alpha,\beta}
t^{-\frac{|\alpha|+|\beta|}{2}}
V(x,y,\sqrt t)^{-1}
\exp\left(
-c_{\alpha,\beta}\frac{d(x,y)^2}{t}
\right)
\end{equation}
for all $t>0$ and $x,y\in\mathbb R^N$.

By \eqref{twosideinequ},
\[
w(B(z,s))
\gtrsim
s^N\prod_{v\in R}
\bigl(|\langle z,v\rangle|+s\bigr)^{k(v)}
\geq
s^N\prod_{v\in R}s^{k(v)}
=
s^\N.
\]
Therefore,
\[
V(x,y,\sqrt t)
=
\max\{w(B(x,\sqrt t)),w(B(y,\sqrt t))\}
\gtrsim
t^{\N/2},
\]
and hence
\[
V(x,y,\sqrt t)^{-1}
\lesssim
t^{-\N/2}.
\]
Substituting this into
\eqref{eq:Dunkl-heat-mixed-lemma}, we obtain
\begin{equation}\label{eq:Dunkl-heat-simple-lemma}
\left|
\partial_x^\alpha\partial_y^\beta h_t(x,y)
\right|
\leq
C_{\alpha,\beta}
t^{-\frac{\N+|\alpha|+|\beta|}{2}}
\exp\left(
-c_{\alpha,\beta}\frac{d(x,y)^2}{t}
\right).
\end{equation}
By functional calculus,
\[
(-\Delta_{\D})^{-1/2}
=
\frac1{\sqrt\pi}
\int_0^\infty
e^{t\Delta_{\D}}\,\frac{dt}{\sqrt t}.
\]
 Consequently, away from the orbit diagonal,
the corresponding integral kernel is
\begin{equation}\label{eq:Dunkl-fractional-kernel-lemma}
K_{\D}(x,y)
=
\frac1{\sqrt\pi}
\int_0^\infty
h_t(x,y)\,\frac{dt}{\sqrt t},
\qquad d(x,y)>0.
\end{equation}
%
%We first verify that the integral on the right-hand side is absolutely
%convergent. By \eqref{eq:Dunkl-heat-simple-lemma} with
%$\alpha=\beta=0$,
%\[
%|h_t(x,y)|
%\lesssim
%t^{-\N/2}
%\exp\left(
%-c\frac{d(x,y)^2}{t}
%\right).
%\]
%Hence, setting $\delta=d(x,y)>0$,
%\[
%\int_0^\infty
%|h_t(x,y)|\,\frac{dt}{\sqrt t}
%\lesssim
%\int_0^\infty
%t^{-\frac{\N+1}{2}}
%e^{-c\delta^2/t}\,dt.
%\]
%With the change of variables
%\[
%s=\frac{\delta^2}{t},
%\qquad
%t=\frac{\delta^2}{s},
%\qquad
%dt=-\delta^2s^{-2}\,ds,
%\]
%we obtain
%\begin{align*}
%\int_0^\infty
%t^{-\frac{\N+1}{2}}
%e^{-c\delta^2/t}\,dt
%&=
%\delta^{1-\N}
%\int_0^\infty
%s^{\frac{\N-3}{2}}e^{-cs}\,ds.
%\end{align*}
%The last integral is finite because $\N>1$. This proves the absolute
%convergence of \eqref{eq:Dunkl-fractional-kernel-lemma}.

%We next justify differentiation under the integral sign. Let
%\[
%\mathcal K
%\Subset
%\{(x,y)\in\mathbb R^N\times\mathbb R^N:d(x,y)>0\}
%\]
%be compact. Since $d$ is continuous,
%\[
%\delta_{\mathcal K}
%:=
%\inf_{(x,y)\in\mathcal K}d(x,y)
%>0.
%\]
%Thus, by \eqref{eq:Dunkl-heat-simple-lemma},
%\[
%\left|
%\partial_x^\alpha\partial_y^\beta h_t(x,y)
%\right|
%\leq
%C_{\alpha,\beta}
%t^{-\frac{\N+|\alpha|+|\beta|}{2}}
%\exp\left(
%-c_{\alpha,\beta}
%\frac{\delta_{\mathcal K}^2}{t}
%\right)
%\]
%uniformly for $(x,y)\in\mathcal K$. After multiplication by
%$t^{-1/2}$, the resulting majorant is integrable on $(0,\infty)$,
%because
%\[
%\N+|\alpha|+|\beta|>1.
%\]
%Therefore dominated convergence permits differentiation under the
%integral in \eqref{eq:Dunkl-fractional-kernel-lemma}. Hence
%\[
%\partial_x^\alpha\partial_y^\beta K_{\D}(x,y)
%=
%\frac1{\sqrt\pi}
%\int_0^\infty
%\partial_x^\alpha\partial_y^\beta h_t(x,y)
%\,\frac{dt}{\sqrt t}
%\]
%whenever $d(x,y)>0$.

Set
\[
m:=|\alpha|+|\beta|.
\]
Using \eqref{eq:Dunkl-heat-simple-lemma}, we obtain
\begin{align*}
\left|
\partial_x^\alpha\partial_y^\beta K_{\D}(x,y)
\right|
&\lesssim
\int_0^\infty
t^{-\frac{\N+m+1}{2}}
\exp\left(
-c\frac{d(x,y)^2}{t}
\right)\,dt.
\end{align*}
Therefore
\[
\left|
\partial_x^\alpha\partial_y^\beta K_{\D}(x,y)
\right|
\lesssim
d(x,y)^{1-\N-m}
=
d(x,y)^{-(\N-1+|\alpha|+|\beta|)},
\]
which proves \eqref{eq:Dunkl-fractional-kernel-derivative-lemma}.

The proof is complete.
\end{proof}

\begin{lemma}\label{lem:rank-one-decomposition-Dunkl}
Assume that
\[
\operatorname{span}(R)=\mathbb R^N,
\qquad
\N>1,
\qquad
2<r<\infty.
\]
Let $K_{\D}(x,y)$ denote the integral kernel of
$(-\Delta_{\D})^{-1/2}$.

Fix a closed Weyl chamber $\Gamma$, and let
\[
\Phi:\mathbb R^N\longrightarrow\mathbb R^N
\]
be the linear isomorphism appearing in
Corollary~\ref{coro:chamber-orthant}. Since
$\operatorname{span}(R)=\mathbb R^N$, we have
\[
\Phi(\Gamma)=\Omega:=\mathbb R_+^N.
\]
Let
\[
\mathbb D_\Gamma
=
\bigcup_{j\in\mathbb Z}\mathbb D_{\Gamma,j}
\]
be the system of dyadic cubes over $\Gamma$ constructed in
Corollary~\ref{Construction over gamma}. Thus, if
$I\in\mathbb D_{\Gamma,j}$, then
\[
I=\Phi^{-1}(Q)
\]
for a unique standard dyadic cube $Q\subset\Omega$ with
$\ell(Q)=2^{-j}$, and we put
\[
\ell(I):=\ell(Q).
\]

For $I=\Phi^{-1}(Q)\in\mathbb D_\Gamma$, let $2Q$ and $3Q$
denote the cubes concentric with $Q$ having side lengths
$2\ell(Q)$ and $3\ell(Q)$, respectively, and set
\[
I^\ast
:=
\Phi^{-1}\bigl((2Q)\cap\Omega\bigr),
\qquad
3_\Phi I
:=
\Phi^{-1}\bigl((3Q)\cap\Omega\bigr).
\]
Then
\[
I\subset I^\ast\subset 3_\Phi I.
\]

There exist nonnegative numbers
$\{c_{p,q}\}_{p,q\in\mathbb Z^N}$, depending only on
$N,R,k$ and the fixed map $\Phi$, such that, for every $M>0$,
\begin{equation}\label{eq:Dunkl-cpq-decay}
c_{p,q}
\leq
C_M(1+|p|+|q|)^{-M},
\qquad
p,q\in\mathbb Z^N.
\end{equation}

For every
$f\in L^r_{\mathrm{loc}}(\mathbb R^N,dw)$ there are functions
\[
e_{p,q,I}^{\sigma,\tau},
\quad
g_{p,q,I}^{\sigma,\tau},
\qquad
p,q\in\mathbb Z^N,\quad
I\in\mathbb D_\Gamma,\quad
\sigma,\tau\in G,
\]
such that, for almost every $(x,y)$ satisfying $d(x,y)>0$,
\begin{equation}\label{eq:Dunkl-Riesz-rank-one}
\begin{split}
f(x)K_{\D}(x,y)
&=
\sum_{\sigma,\tau\in G}
\sum_{p,q\in\mathbb Z^N}
c_{p,q}
\sum_{I\in\mathbb D_\Gamma}
\ell(I)^{-(\N-1)}
w(I)^{1-\frac1r}
\|f\|_{L^r(\sigma I,dw)}
e_{p,q,I}^{\sigma,\tau}(x)
g_{p,q,I}^{\sigma,\tau}(y).
\end{split}
\end{equation}
The functions in \eqref{eq:Dunkl-Riesz-rank-one} satisfy
\begin{equation}\label{eq:Dunkl-support-e}
\operatorname{supp}e_{p,q,I}^{\sigma,\tau}
\subset \sigma I
\subset \mathcal O(3_\Phi I)
\end{equation}
and
\begin{equation}\label{eq:Dunkl-support-g}
\operatorname{supp}g_{p,q,I}^{\sigma,\tau}
\subset \tau I^\ast
\subset \tau(3_\Phi I)
\subset \mathcal O(3_\Phi I),
\end{equation}
\begin{equation}\label{eq:Dunkl-normalization-e}
\sup_{\substack{\sigma,\tau\in G\\
p,q\in\mathbb Z^N,\ I\in\mathbb D_\Gamma}}
w(I)^{\frac12-\frac1r}
\|e_{p,q,I}^{\sigma,\tau}\|_{L^r(\mathbb R^N,dw)}
\leq 1,
\end{equation}
\begin{equation}\label{eq:Dunkl-normalization-g}
\sup_{\substack{\sigma,\tau\in G\\
p,q\in\mathbb Z^N,\ I\in\mathbb D_\Gamma}}
w(I)^{\frac12-\frac1r}
\|g_{p,q,I}^{\sigma,\tau}\|_{L^r(\mathbb R^N,dw)}
\leq 1.
\end{equation}
\end{lemma}

\begin{proof}
We divide the proof into several steps.


Fix $\sigma,\tau\in G$, and suppose that
$
x\in\sigma(\Gamma)^\circ,
$ and $
y\in\tau(\Gamma)^\circ.
$
Set
$
x_\sigma:=\sigma^{-1}x\in\Gamma^\circ,$ and $
y_\tau:=\tau^{-1}y\in\Gamma^\circ,
$
and
$
u:=\Phi(x_\sigma),$ and
$
v:=\Phi(y_\tau).
$

Since $x_\sigma,y_\tau$ belong to the same closed Weyl chamber,
\eqref{distanceequa} gives
\[
d(x,y)
=
d(x_\sigma,y_\tau)
=
\|x_\sigma-y_\tau\|.
\]
Together with \eqref{eq:biLip-N}, this yields
\begin{equation}\label{eq:d-Phi-comparable}
C_\Phi^{-1}\|u-v\|
\leq
d(x,y)
\leq
c_\Phi^{-1}\|u-v\|.
\end{equation}

Choose $\rho\in C_c^\infty(\mathbb R^N)$ such that
\[
0\leq\rho\leq1,
\qquad
\rho(z)=1
\quad\text{if }|z|_\infty\leq\frac14,
\]
and
\[
\rho(z)=0
\quad\text{if }|z|_\infty\geq\frac12.
\]
Define
\[
\psi(z):=\rho(z)-\rho(2z).
\]
Then
\begin{equation}\label{eq:psi-support-Dunkl}
\operatorname{supp}\psi
\subset
\left\{
z:
\frac18\leq |z|_\infty\leq\frac12
\right\},
\end{equation}
and, for every $z\neq0$,
\begin{equation}\label{eq:psi-partition-Dunkl}
\sum_{j\in\mathbb Z}\psi(2^jz)=1.
\end{equation}
Indeed,
\[
\sum_{j=-m}^{n}\psi(2^jz)
=
\rho(2^{-m}z)-\rho(2^{n+1}z),
\]
and letting $m,n\to\infty$ proves
\eqref{eq:psi-partition-Dunkl}.

For $j\in\mathbb Z$, and 
for almost every $u\in\Omega$, there is a unique
$Q\in\mathbb D_{\Omega,j}^0$ containing $u$. If
\[
I:=\Phi^{-1}(Q)\in\mathbb D_{\Gamma,j},
\]
then $\ell(I)=2^{-j}$. Therefore,
for $x\in\sigma(\Gamma)^\circ$ and
$y\in\tau(\Gamma)^\circ$ with $d(x,y)>0$,
\begin{equation}\label{eq:K-shell-decomp}
K_{\D}(x,y)
=
\sum_{j\in\mathbb Z}
\sum_{I\in\mathbb D_{\Gamma,j}}
\chi_{\sigma I}(x)
K_{\D}(x,y)
\psi\left(
\frac{
\Phi(\sigma^{-1}x)-\Phi(\tau^{-1}y)
}{\ell(I)}
\right).
\end{equation}
For each fixed $(x,y)$ with $d(x,y)>0$, only finitely many
$j$ contribute to the sum, by
\eqref{eq:psi-support-Dunkl}.

Write
\[
Q=\prod_{i=1}^N[m_i2^{-j},(m_i+1)2^{-j}),
\qquad
m_i\in\mathbb Z_+,
\]
and denote its center by $c_Q$. Suppose that
$
u\in Q
$
and
\[
\psi\left(\frac{u-v}{\ell}\right)\neq0.
\]
Then
$
|u-v|_\infty\leq\frac{\ell}{2}.
$
Since
$
|u-c_Q|_\infty\leq\frac{\ell}{2},
$
we obtain
$
|v-c_Q|_\infty\leq\ell.
$
Hence
$
v\in(2Q)\cap\Omega.
$
Equivalently,
\begin{equation}\label{eq:y-localization-Dunkl}
y\in\tau I^\ast.
\end{equation}
In particular,
\[
\tau I^\ast
\subset\tau(3_\Phi I)
\subset\mathcal O(3_\Phi I).
\]
Moreover, by the doubling property, we see that
\begin{equation}\label{eq:wIstar-wI}
w(I)
\leq
w(I^\ast)
\leq
C\,w(I)
\end{equation}
uniformly in $I\in\mathbb D_\Gamma$.

Fix $\sigma,\tau\in G$ and
$I=\Phi^{-1}(Q)\in\mathbb D_\Gamma$. Write
$
\ell:=\ell(I)
$
and let $c_Q$ be the center of $Q$.

For
\[
a=\frac{u-c_Q}{\ell},
\qquad
b=\frac{v-c_Q}{\ell},
\]
the condition $u\in Q$ means
\[
a\in A:=[-1/2,1/2]^N.
\]
Furthermore, the condition
$
v\in(2Q)\cap\Omega
$
is equivalent to
\[
b\in B_Q
:=
\left\{
b\in[-1,1]^N:
c_Q+\ell b\in\Omega
\right\}.
\]
Since
\[
(c_Q)_i=(m_i+1/2)\ell,
\qquad m_i\in\mathbb Z_+,
\]
we have
\[
B_Q
=
\prod_{i=1}^N
\left[
\max\{-1,-m_i-1/2\},1
\right].
\]
Thus the left endpoint of the $i$-th interval is equal to
$-1/2$ if $m_i=0$ and is equal to $-1$ if $m_i\geq1$.
Consequently, the rectangles
$
A\times B_Q
$
belong to a finite family containing at most $2^N$ different
rectangles.

We claim that there exists a number $\delta>0$, depending only
on $N,R,k$ and $\Phi$, with the following property. For every
$Q$ one can choose
$
\chi_Q\in C_c^\infty((-2,2)^N\times(-2,2)^N)
$
such that
\[
\chi_Q=1
\quad\text{on }A\times B_Q,
\]
\[
\operatorname{supp}\chi_Q
\subset
\left\{
(a,b):
\operatorname{dist}\bigl((a,b),A\times B_Q\bigr)<\delta
\right\},
\]
and, for every integer $M\geq0$,
\begin{equation}\label{eq:chi-uniform}
\sup_Q\|\chi_Q\|_{C^M}
<\infty.
\end{equation}
Indeed, there are only finitely many possible rectangles
$A\times B_Q$, and hence fixed smooth cutoff functions can be
chosen for these finitely many rectangles.

We choose $\delta>0$ sufficiently small so that the following
additional property holds:
\begin{equation}\label{eq:orbit-separation-neighborhood}
d\left(
\sigma\Phi^{-1}(c_Q+\ell a),
\tau\Phi^{-1}(c_Q+\ell b)
\right)
\geq c_0\ell
\end{equation}
whenever
$
(a,b)\in\operatorname{supp}\chi_Q
$
and
$
a-b\in\operatorname{supp}\psi.
$
Let us verify this assertion.

Take $(a,b)\in\operatorname{supp}\chi_Q$. There exists
$(a_0,b_0)\in A\times B_Q$ such that
\[
|a-a_0|+|b-b_0|\leq C_N\delta.
\]
If $a-b\in\operatorname{supp}\psi$, then
$
|a-b|_\infty\geq\frac18.
$
By decreasing $\delta$ if necessary, we may assume that
$
|a_0-b_0|_\infty\geq\frac1{16}.
$
Set
\[
x_0'
=
\Phi^{-1}(c_Q+\ell a_0),
\qquad
y_0'
=
\Phi^{-1}(c_Q+\ell b_0).
\]
Since
$
c_Q+\ell a_0,\ c_Q+\ell b_0\in\Omega,
$
we have $x_0',y_0'\in\Gamma$. Hence
\eqref{distanceequa} and \eqref{eq:biLip-N} imply
\begin{align*}
d(\sigma x_0',\tau y_0')
&=
d(x_0',y_0')
\\
&=
\|x_0'-y_0'\|
\\
&\geq
C_\Phi^{-1}
\ell|a_0-b_0|
\\
&\geq
c_1\ell.
\end{align*}
On the other hand, note that the orbit distance satisfies
\begin{equation}\label{eq:d-Lipschitz}
|d(x,y)-d(x',y')|
\leq
\|x-x'\|+\|y-y'\|.
\end{equation}

Since $\Phi^{-1}$ is linear and bounded,
\[
\left\|
\Phi^{-1}(c_Q+\ell a)
-
\Phi^{-1}(c_Q+\ell a_0)
\right\|
\lesssim\ell\delta,
\]
and similarly for $b,b_0$. Thus
\eqref{eq:d-Lipschitz} gives
\[
d\left(
\sigma\Phi^{-1}(c_Q+\ell a),
\tau\Phi^{-1}(c_Q+\ell b)
\right)
\geq
c_1\ell-C\delta\ell.
\]
Choosing $\delta$ so small that
$C\delta\leq c_1/2$ proves
\eqref{eq:orbit-separation-neighborhood}.

Define
\begin{equation}\label{eq:def-H-Dunkl}
\begin{split}
H_{\sigma,\tau,I}(a,b)
:={}&
\ell^{\N-1}
\chi_Q(a,b)\psi(a-b)
K_{\D}\left(
\sigma\Phi^{-1}(c_Q+\ell a),
\tau\Phi^{-1}(c_Q+\ell b)
\right).
\end{split}
\end{equation}
On the support of
$\chi_Q(a,b)\psi(a-b)$,
\eqref{eq:orbit-separation-neighborhood} gives
$$d\left(
\sigma\Phi^{-1}(c_Q+\ell a),
\tau\Phi^{-1}(c_Q+\ell b)
\right)\gtrsim\ell.$$

We claim that, for every integer $M\geq0$,
\begin{equation}\label{eq:H-CM-uniform}
\sup_{\sigma,\tau\in G}
\sup_{I\in\mathbb D_\Gamma}
\|H_{\sigma,\tau,I}\|_{C^M(\mathbb R^{2N})}
\leq C_M.
\end{equation}
Indeed, consider a derivative of total order $m$ in $(a,b)$.
Since $\Phi^{-1}$, $\sigma$ and $\tau$ are fixed linear maps,
the chain rule give
\begin{align*}
&\ell^{\N-1}\ell^m
\left|
\partial_x^\alpha\partial_y^\beta
K_{\D}\left(
\sigma\Phi^{-1}(c_Q+\ell a),
\tau\Phi^{-1}(c_Q+\ell b)
\right)
\right|
\\
&\qquad\lesssim
\ell^{\N-1+m}
(c_0\ell)^{-(\N-1+m)}
\lesssim1,
\end{align*}
where $|\alpha|+|\beta|=m$. Combining this with
\eqref{eq:chi-uniform} and the fixed smoothness of $\psi$
proves \eqref{eq:H-CM-uniform}.

Since $H_{\sigma,\tau,I}$ is supported in the interior of
$
[-2,2]^N\times[-2,2]^N,
$
its $4$-periodic extension in both variables is smooth. Hence
it has the Fourier expansion
\begin{equation}\label{eq:Dunkl-local-Fourier}
H_{\sigma,\tau,I}(a,b)
=
\sum_{p,q\in\mathbb Z^N}
b_{p,q,I}^{\sigma,\tau}
e^{\frac{\pi i}{2}p\cdot a}
e^{\frac{\pi i}{2}q\cdot b},
\end{equation}
where
\begin{equation}\label{eq:b-coefficients}
\begin{split}
b_{p,q,I}^{\sigma,\tau}
={}&
4^{-2N}
\int_{[-2,2]^N}
\int_{[-2,2]^N}
H_{\sigma,\tau,I}(a,b)
e^{-\frac{\pi i}{2}p\cdot a}
e^{-\frac{\pi i}{2}q\cdot b}
\,da\,db.
\end{split}
\end{equation}
Repeated integration by parts, together with
\eqref{eq:H-CM-uniform}, gives, for every $M>0$,
\begin{equation}\label{eq:b-rapid}
\sup_{\sigma,\tau\in G}
\sup_{I\in\mathbb D_\Gamma}
|b_{p,q,I}^{\sigma,\tau}|
\leq
C_M(1+|p|+|q|)^{-M}.
\end{equation}

For $a\in A$ and $b\in B_Q$, we have $\chi_Q(a,b)=1$.
Consequently, \eqref{eq:Dunkl-local-Fourier} implies
\begin{equation}\label{eq:local-K-Fourier}
\begin{split}
&K_{\D}(x,y)
\psi\left(
\frac{\Phi(\sigma^{-1}x)-\Phi(\tau^{-1}y)}
{\ell(I)}
\right)
\\
&\quad=
\ell(I)^{-(\N-1)}
\sum_{p,q\in\mathbb Z^N}
b_{p,q,I}^{\sigma,\tau}
\exp\left(
\frac{\pi i}{2\ell(I)}
p\cdot
\bigl(\Phi(\sigma^{-1}x)-c_Q\bigr)
\right)
\exp\left(
\frac{\pi i}{2\ell(I)}
q\cdot
\bigl(\Phi(\tau^{-1}y)-c_Q\bigr)
\right)
\end{split}
\end{equation}
whenever
$
x\in\sigma I,$ and $
y\in\tau I^\ast.
$

The coefficients in \eqref{eq:local-K-Fourier} depend on
$\sigma,\tau$ and $I$. We now absorb this harmless dependence
into the normalized rank-one factors.

Define
\begin{equation}\label{eq:def-cpq}
c_{p,q}
:=
\sup_{\substack{\sigma,\tau\in G\\I\in\mathbb D_\Gamma}}
\left\{
|b_{p,q,I}^{\sigma,\tau}|
\left(\frac{w(I^\ast)}{w(I)}\right)^{1/2}
\right\}.
\end{equation}
By \eqref{eq:wIstar-wI} and \eqref{eq:b-rapid},
for every $M>0$,
\[
c_{p,q}
\leq
C_M(1+|p|+|q|)^{-M}.
\]
This proves \eqref{eq:Dunkl-cpq-decay}.

For each $\sigma,\tau,p,q,I$, define
\begin{equation}\label{eq:def-a-pq}
a_{p,q,I}^{\sigma,\tau}
:=
\begin{cases}
\displaystyle
\frac{b_{p,q,I}^{\sigma,\tau}}{c_{p,q}}
\left(\frac{w(I^\ast)}{w(I)}\right)^{1/2},
& c_{p,q}>0,\\[3mm]
0,&c_{p,q}=0.
\end{cases}
\end{equation}
By definition,
\begin{equation}\label{eq:a-bound}
|a_{p,q,I}^{\sigma,\tau}|\leq1.
\end{equation}
Moreover, if $c_{p,q}=0$, then necessarily
$b_{p,q,I}^{\sigma,\tau}=0$ for every $\sigma,\tau,I$.
Thus, in all cases,
\begin{equation}\label{eq:b-factorization}
b_{p,q,I}^{\sigma,\tau}
=
c_{p,q}
a_{p,q,I}^{\sigma,\tau}
\left(\frac{w(I)}{w(I^\ast)}\right)^{1/2}.
\end{equation}


Fix $\sigma,\tau\in G$ and
$I\in\mathbb D_\Gamma$.
If $\|f\|_{L^r(\sigma I,dw)}>0$, define
\begin{equation}\label{eq:def-e-Dunkl}
\begin{split}
e_{p,q,I}^{\sigma,\tau}(x)
:={}&
a_{p,q,I}^{\sigma,\tau}
w(I)^{\frac1r-\frac12}
\frac{f(x)\chi_{\sigma I}(x)}
{\|f\|_{L^r(\sigma I,dw)}}
\exp\left(
\frac{\pi i}{2\ell(I)}
p\cdot
\bigl(\Phi(\sigma^{-1}x)-c_Q\bigr)
\right).
\end{split}
\end{equation}
If $\|f\|_{L^r(\sigma I,dw)}=0$, set
$
e_{p,q,I}^{\sigma,\tau}:=0.
$
Also define
\begin{equation}\label{eq:def-g-Dunkl}
\begin{split}
g_{p,q,I}^{\sigma,\tau}(y)
:={}&
w(I^\ast)^{-1/2}
\chi_{\tau I^\ast}(y)
\exp\left(
\frac{\pi i}{2\ell(I)}
q\cdot
\bigl(\Phi(\tau^{-1}y)-c_Q\bigr)
\right).
\end{split}
\end{equation}
The support properties
\eqref{eq:Dunkl-support-e}--\eqref{eq:Dunkl-support-g}
are immediate.

Since the Dunkl measure is $G$-invariant,
\[
w(\sigma I)=w(I),
\qquad
w(\tau I^\ast)=w(I^\ast).
\]
From \eqref{eq:a-bound} and
\eqref{eq:def-e-Dunkl}, we obtain
\begin{align*}
\|e_{p,q,I}^{\sigma,\tau}\|_{L^r(dw)}
&\leq
w(I)^{\frac1r-\frac12}
\frac{
\|f\chi_{\sigma I}\|_{L^r(dw)}
}{
\|f\|_{L^r(\sigma I,dw)}
}
=
w(I)^{\frac1r-\frac12}.
\end{align*}
Therefore,
\[
w(I)^{\frac12-\frac1r}
\|e_{p,q,I}^{\sigma,\tau}\|_{L^r(dw)}
\leq1,
\]
which proves \eqref{eq:Dunkl-normalization-e}.

Similarly,
\[
\|g_{p,q,I}^{\sigma,\tau}\|_{L^r(dw)}
=
w(I^\ast)^{\frac1r-\frac12}.
\]
Hence
\begin{align*}
&w(I)^{\frac12-\frac1r}
\|g_{p,q,I}^{\sigma,\tau}\|_{L^r(dw)}=
\left(
\frac{w(I)}{w(I^\ast)}
\right)^{\frac12-\frac1r}
\leq1.
\end{align*}

This proves \eqref{eq:Dunkl-normalization-g}.


Assume first that
$
\|f\|_{L^r(\sigma I,dw)}>0.
$
Combining \eqref{eq:def-e-Dunkl},
\eqref{eq:def-g-Dunkl} and
\eqref{eq:b-factorization}, we obtain, for
$x\in\sigma I$ and $y\in\tau I^\ast$,
\begin{align}
&c_{p,q}
w(I)^{1-\frac1r}
\|f\|_{L^r(\sigma I,dw)}
e_{p,q,I}^{\sigma,\tau}(x)
g_{p,q,I}^{\sigma,\tau}(y)
\nonumber\\
&\quad=
b_{p,q,I}^{\sigma,\tau}
f(x)
\exp\left(
\frac{\pi i}{2\ell(I)}
p\cdot
\bigl(\Phi(\sigma^{-1}x)-c_Q\bigr)
\right)
\exp\left(
\frac{\pi i}{2\ell(I)}
q\cdot
\bigl(\Phi(\tau^{-1}y)-c_Q\bigr)
\right).
\label{eq:rank-one-exact-identity}
\end{align}
Indeed, the powers of the measures cancel according to
\[
w(I)^{1-\frac1r}
w(I)^{\frac1r-\frac12}
w(I^\ast)^{-1/2}
\left(\frac{w(I^\ast)}{w(I)}\right)^{1/2}
=1.
\]

Multiplying \eqref{eq:local-K-Fourier} by $f(x)$ and using
\eqref{eq:rank-one-exact-identity}, we get
\begin{align}
&f(x)K_{\D}(x,y)
\psi\left(
\frac{
\Phi(\sigma^{-1}x)-\Phi(\tau^{-1}y)
}{\ell(I)}
\right)
\nonumber\\
&\quad=
\sum_{p,q\in\mathbb Z^N}
c_{p,q}
\ell(I)^{-(\N-1)}
w(I)^{1-\frac1r}
\|f\|_{L^r(\sigma I,dw)}
e_{p,q,I}^{\sigma,\tau}(x)
g_{p,q,I}^{\sigma,\tau}(y).
\label{eq:one-piece-reconstruction}
\end{align}
If $\|f\|_{L^r(\sigma I,dw)}=0$, then $f=0$ almost everywhere on
$\sigma I$, and hence \eqref{eq:one-piece-reconstruction}
continues to hold for almost every $x\in\sigma I$.

Moreover, if $x\in\sigma I$ but
$y\notin\tau I^\ast$, then
\eqref{eq:y-localization-Dunkl} shows that
\[
\psi\left(
\frac{
\Phi(\sigma^{-1}x)-\Phi(\tau^{-1}y)
}{\ell(I)}
\right)=0,
\]
while
$
g_{p,q,I}^{\sigma,\tau}(y)=0.
$
Thus \eqref{eq:one-piece-reconstruction} is valid for almost
every
$
(x,y)\in\sigma I\times\tau(\Gamma).
$

We now sum \eqref{eq:one-piece-reconstruction} over
$I\in\mathbb D_{\Gamma,j}$ and then over $j\in\mathbb Z$.
Using \eqref{eq:K-shell-decomp}, we obtain, for almost every
$
x\in\sigma(\Gamma)^\circ,
y\in\tau(\Gamma)^\circ,
d(x,y)>0,
$
that
\begin{align*}
f(x)K_{\D}(x,y)
&=
\sum_{p,q\in\mathbb Z^N}
c_{p,q}
\sum_{I\in\mathbb D_\Gamma}
\ell(I)^{-(\N-1)}
w(I)^{1-\frac1r}
\|f\|_{L^r(\sigma I,dw)}
e_{p,q,I}^{\sigma,\tau}(x)
g_{p,q,I}^{\sigma,\tau}(y).
\end{align*}

Finally, the interiors of the Weyl chambers
$\{\sigma(\Gamma):\sigma\in G\}$ are pairwise disjoint and
cover $\mathbb R^N$ up to the union of the reflecting
hyperplanes, which is a $dw$-null set. Summing the last
identity over $\sigma,\tau\in G$ gives exactly
\eqref{eq:Dunkl-Riesz-rank-one}.

This completes the proof.
\end{proof}


\begin{lemma}\label{lem:counterexample-weighted-Cwikel}
Suppose that
\[
dw(x)=w(x)\,dx,
\qquad
w(x)=\prod_{v\in R}|\langle x,v\rangle|^{k(v)},
\]
is a nontrivial Dunkl weight, that is,
\[
\kappa:=\sum_{v\in R}k(v)>0,
\]
and let $\Gamma$ be a Weyl chamber. Then the following assertion is
false in general: if an integral operator $V$ on
$L^2(\Gamma,dw)$ has a kernel $K(x,y)$ satisfying
\[
\left|
\partial_x^\alpha\partial_y^\beta K(x,y)
\right|
\leq
C_{\alpha,\beta}
|x-y|^{-(N-1+|\alpha|+|\beta|)},
\qquad x,y\in\Gamma,\quad x\neq y,
\]
then
\[
\|M_fV\|_{S^{N,\infty}(L^2(\Gamma,dw))}
\lesssim
\|f\|_{L^N(\Gamma,dx)}
\]
for every $f\in L^N(\Gamma,dx)$, with an implicit constant depending
only on the constants in the kernel estimates and on the fixed Dunkl
data.

More precisely, there exist kernels $K_R$, with the constants in the
above derivative estimates independent of $R$, and functions
$f_R\in C_c^\infty(\Gamma)$ such that
\[
\sup_{R\gg1}\|f_R\|_{L^N(\Gamma,dx)}<\infty,
\]
whereas
\[
\|M_{f_R}V_R\|_{S^{N,\infty}(L^2(\Gamma,dw))}
\longrightarrow\infty
\qquad\text{as }R\to\infty.
\]
\end{lemma}

\begin{proof}
Choose $\xi\in\Gamma^\circ$. Since $\xi$ lies in the interior of a
Weyl chamber,
\[
\langle \xi,v\rangle\neq0
\]
for every $v\in R$. Hence, for every fixed $C>0$,
\begin{equation}\label{eq:weight-along-chamber-ray}
w(R\xi+z)\asymp R^\kappa,
\qquad |z|\leq C,
\end{equation}
uniformly for all sufficiently large $R$. Indeed,
\[
\begin{aligned}
w(R\xi+z)
&=
\prod_{v\in R}
\left|
R\langle\xi,v\rangle+\langle z,v\rangle
\right|^{k(v)}
\\
&=
R^\kappa
\prod_{v\in R}
\left|
\langle\xi,v\rangle
+
R^{-1}\langle z,v\rangle
\right|^{k(v)},
\end{aligned}
\]
Since $\xi\in\Gamma^\circ$ and the root system $R$ is finite,
\[
m_\xi:=\min_{v\in R}|\langle\xi,v\rangle|>0.
\]
Moreover, for $|z|\leq C$,
\[
|R^{-1}\langle z,v\rangle|
\leq R^{-1}C\max_{u\in R}|u|.
\]
Thus, for all sufficiently large $R$, uniformly in $|z|\leq C$ and
$v\in R$,
\[
\frac12|\langle\xi,v\rangle|
\leq
\left|\langle\xi,v\rangle+R^{-1}\langle z,v\rangle\right|
\leq
\frac32|\langle\xi,v\rangle|.
\]
Since $k(v)\geq0$ and $R$ is finite, it follows that
\[
0<c_\xi
\leq
\prod_{v\in R}
\left|\langle\xi,v\rangle+R^{-1}\langle z,v\rangle\right|^{k(v)}
\leq C_\xi<\infty
\]
uniformly for $|z|\leq C$ and all sufficiently large $R$.

Fix nonzero functions
\[
\phi,\psi\in C_c^\infty(B(0,\varepsilon)),
\]
where $\varepsilon>0$ is sufficiently small, and choose a fixed vector
$a\in\mathbb R^N$ such that
\[
|a|>4\varepsilon.
\]
Set
\[
x_R:=R\xi,
\qquad
y_R:=R\xi+a.
\]
Since $\xi\in\Gamma^\circ$ and $a$ is fixed, for all sufficiently
large $R$,
\[
B(x_R,\varepsilon)\cup B(y_R,\varepsilon)
\subset\Gamma^\circ.
\]

Define
\[
K_R(x,y)
:=
\phi(x-x_R)\overline{\psi(y-y_R)},
\qquad x,y\in\Gamma,
\]
and let $V_R$ be the corresponding integral operator on
$L^2(\Gamma,dw)$:
\[
V_Ru(x)
=
\int_\Gamma K_R(x,y)u(y)\,dw(y).
\]
For every pair of multi-indices $\alpha,\beta$,
\[
\partial_x^\alpha\partial_y^\beta K_R(x,y)
=
(\partial^\alpha\phi)(x-x_R)
\overline{(\partial^\beta\psi)(y-y_R)}.
\]
If this expression is nonzero, then
\[
|x-x_R|<\varepsilon,
\qquad
|y-y_R|<\varepsilon,
\]
and therefore
\[
|a|-2\varepsilon
\leq |x-y|
\leq |a|+2\varepsilon.
\]
Thus $|x-y|\asymp1$ on the support of every derivative of $K_R$, with
constants independent of $R$. It follows that
\[
\left|
\partial_x^\alpha\partial_y^\beta K_R(x,y)
\right|
\leq
C_{\alpha,\beta}
|x-y|^{-(N-1+|\alpha|+|\beta|)}
\]
for all $x\neq y$, where $C_{\alpha,\beta}$ is independent of $R$.

Put
\[
\phi_R(x):=\phi(x-x_R),
\qquad
\psi_R(y):=\psi(y-y_R).
\]
Then
\[
V_Ru
=
\phi_R\,\langle u,\psi_R\rangle_{L^2(dw)},
\]
so $V_R$ is a rank-one operator. Its unique nonzero singular value is
\[
s_1(V_R)
=
\|\phi_R\|_{L^2(dw)}
\|\psi_R\|_{L^2(dw)}.
\]
By \eqref{eq:weight-along-chamber-ray},
\[
\begin{aligned}
\|\phi_R\|_{L^2(dw)}^2
&=
\int_\Gamma
|\phi(x-x_R)|^2w(x)\,dx
\\
&=
\int_{\mathbb R^N}
|\phi(z)|^2w(R\xi+z)\,dz
\\
&\asymp
R^\kappa
\int_{\mathbb R^N}|\phi(z)|^2\,dz,
\end{aligned}
\]
and hence
\[
\|\phi_R\|_{L^2(dw)}
\asymp R^{\kappa/2}.
\]
Likewise,
\[
\|\psi_R\|_{L^2(dw)}
\asymp R^{\kappa/2}.
\]
Consequently,
\begin{equation}\label{eq:rank-one-growth-counterexample}
s_1(V_R)\asymp R^\kappa.
\end{equation}

Choose $\eta\in C_c^\infty(\mathbb R^N)$ such that
\[
\eta=1
\quad\text{on }\operatorname{supp}\phi,
\]
and define
\[
f_R(x):=\eta(x-x_R).
\]
For all sufficiently large $R$,
\[
\operatorname{supp}f_R\subset\Gamma^\circ,
\]
and by translation invariance of Lebesgue measure,
\[
\|f_R\|_{L^N(\Gamma,dx)}
=
\|\eta\|_{L^N(\mathbb R^N)}.
\]
In particular,
\[
\sup_{R\gg1}
\|f_R\|_{L^N(\Gamma,dx)}
<\infty.
\]

Since $f_R=1$ on $\operatorname{supp}\phi_R$, we have
\[
M_{f_R}V_R=V_R.
\]
As $V_R$ has rank one,
\[
\|V_R\|_{S^{N,\infty}(L^2(\Gamma,dw))}
=
s_1(V_R).
\]
Therefore, by \eqref{eq:rank-one-growth-counterexample},
\[
\|M_{f_R}V_R\|_{S^{N,\infty}(L^2(\Gamma,dw))}
\asymp R^\kappa
\longrightarrow\infty,
\]
whereas
\[
\|f_R\|_{L^N(\Gamma,dx)}
=
\|\eta\|_{L^N(\mathbb R^N)}
\]
is independent of $R$. This proves the assertion.
\end{proof}

\bibliographystyle{plain}
\bibliography{references}

\end{document}
